% Copyright (C) 2025-2026  TOS Network.
% Licensed under the GNU General Public License v3.0.
\documentclass[11pt,oneside]{article}

\usepackage[T1]{fontenc}
\usepackage[english]{babel}
\usepackage{amsmath,amssymb}
\usepackage{array}
\usepackage{booktabs}
\usepackage{enumitem}
\usepackage{geometry}
\usepackage{hyperref}
\usepackage{longtable}
\usepackage{microtype}
\usepackage{xcolor}

\geometry{margin=1in}
\emergencystretch=2em
\hypersetup{
  colorlinks=true,
  linkcolor=blue!55!black,
  urlcolor=blue!55!black,
  pdftitle={TOS Network: The Open System for the Agentic Internet},
  pdfauthor={TOS Network}
}
\setlist{nosep}
\setcounter{tocdepth}{2}

\newcommand{\tos}{\textsc{TOS}}
\newcommand{\nanotos}{\textit{nanotomi}}
\newcommand{\implemented}{\textbf{Implemented}}
\newcommand{\experimental}{\textbf{Experimental}}
\newcommand{\planned}{\textbf{Planned}}
\newcommand{\partialstatus}{\textbf{Partial}}
\newcommand{\opid}[1]{\texttt{#1}}

\newcolumntype{L}[1]{>{\raggedright\arraybackslash}p{#1}}

\makeatletter
\g@addto@macro\UrlBreaks{\do\-}
\let\tos@makefntext\@makefntext
\renewcommand\@makefntext[1]{\tos@makefntext{\raggedright #1}}
\makeatother

\title{
  \textbf{TOS Network}\\[0.4em]
  \Large The Open System for the Agentic Internet\\[0.25em]
  \large Foundational Infrastructure for the Agentic Internet\\[0.2em]
  \normalsize Mission: Connect Every AI Agent
}
\author{TOS Network}
\date{2026 Agentic Internet Edition}

\begin{document}
\maketitle

\begin{abstract}
AI agents are moving from assistants that recommend actions to autonomous participants that
discover counterparties, negotiate terms, invoke services, spend bounded budgets and
coordinate across organizational boundaries. Existing APIs, agent interfaces and payment
rails are necessary, but they do not by themselves establish who an agent represents, what it
may do, what both parties accepted, how execution is bounded, what evidence must be returned
or how failure is resolved.

TOS Network is being built as the \textbf{foundational infrastructure---The Open System---for
the Agentic Internet}. Its mission is \textbf{Connect Every AI Agent}. The common lifecycle is
\textbf{Identity $\rightarrow$ Discovery $\rightarrow$ Intent $\rightarrow$ Negotiation
$\rightarrow$ Agreement $\rightarrow$ Bounded Execution $\rightarrow$ Evidence/Receipt
$\rightarrow$ Settlement}. An Intent expresses a revocable objective; it does not by itself
create an obligation. An Agreement records the exact terms accepted by the relevant parties.
Only then is work delegated to a bounded execution path.

Agentic Service Operations are execution-layer primitives within this larger infrastructure,
not the definition of TOS as a whole. Each versioned operation binds an input and output
contract to capability scope, replay domain, resource envelope, maximum budget, metering,
receipt semantics and settlement policy. Communication, model and tool calls, storage,
events, software work, booking, commerce and physical-edge work can use the same accountable
envelope while retaining domain-specific execution and safety rules.

TOS reuses HTTP, MCP, A2A, OpenAPI and compatible discovery systems rather than replacing
them. Models may interpret user goals, but deterministic policy and explicit authority decide
what can be accepted and executed. Providers retain control over admission, hardware, models,
data, pricing and local safety. Private payloads and workloads remain off-chain; TOS Core is
used for shared identity, commitments, disputes and final settlement where independent
parties need a tamper-resistant source of truth.

The implemented foundation, TOS Core, is an actor-model blockchain with native TVM execution,
value-bearing asynchronous messages, masterchain-rooted consensus, dynamic sharding and
ADNL/DHT/RLDP networking. The open Intent, Agreement and Agentic Operation protocols, SDKs,
conformance suites and most service profiles are partial or planned. Artificial Intelligence
Proof of Work (AIPoW) is feature-gated off by default: token incentives must not substitute
for product-market fit, organic paid demand or sustainable unit economics.
\end{abstract}

\section*{Document Status}

This paper describes both released foundations and architectural direction. ``Implemented''
means the referenced repository contains the relevant node, contract, tooling or test
foundation. ``Partial'' means useful primitives or one profile-specific implementation exist,
but the generic interoperable product surface is incomplete. ``Planned'' identifies work that
still requires specification, implementation, conformance testing and independent review.

\tableofcontents
\newpage

\section{Executive Summary}

The Internet connected computers. TOS is being built to connect autonomous agents.

TOS Network is the proposed foundational infrastructure for the Agentic Internet: an open
system through which agents operated by different people, organizations and runtimes can find
one another, communicate, form precise commitments, act within delegated limits, exchange
evidence and settle outcomes. Its mission is:

\begin{center}
\textbf{Connect Every AI Agent.}
\end{center}

An API may describe how to format and transport a request, yet it usually does not provide a
portable answer to the questions autonomous counterparties must resolve before acting:

\begin{itemize}
  \item who is making the request and under whose authority;
  \item how a goal becomes a proposal and then a mutually accepted Agreement;
  \item which exact operation and revision will execute that Agreement;
  \item what resources the operation may consume;
  \item how much the caller is willing and permitted to spend;
  \item how replay, expiry, cancellation and retries are bounded;
  \item what evidence and receipt the provider must return; and
  \item how payment, refund, dispute and final settlement occur.
\end{itemize}

TOS Network is being built to standardize this missing infrastructure. The primary lifecycle
is not a payment call or an economic opcode in isolation:

\begin{center}
\texttt{Identity -> Discovery -> Intent -> Negotiation -> Agreement}\\
\texttt{-> Bounded Execution -> Evidence/Receipt -> Settlement}
\end{center}

Identity establishes the accountable principal and its delegated authority. Discovery returns
candidates, not permission. Intent expresses a revocable objective. Negotiation refines price,
scope, timing, evidence and recourse. Agreement is the exact accepted commitment. Agentic
Service Operations then provide typed, metered execution primitives. Evidence and Receipts
support verification, reputation, dispute and final settlement.

\begin{center}
\fbox{\parbox{0.90\textwidth}{
\centering
\textbf{Every action is typed. Every authority is bounded.}\\
\textbf{Every resource is metered. Every result is receipted.}\\
\textbf{Every cost is accountable.}
}}
\end{center}

The architecture has five interoperable planes:

\begin{enumerate}
  \item \textbf{Trust and Settlement Plane.} TOS Core consensus, TVM, actor accounts,
        commitments, assets, payments, disputes and finality.
  \item \textbf{Identity and Operation Plane.} Agents, capabilities, Intent, Agreement,
        versioned operations, authorization, metering, evidence and Receipts.
  \item \textbf{Propagation and Discovery Plane.} Direct links, ADNL/DHT/RLDP, TOS DNS,
        TOS Sites, Messenger, Mailbox, gateways, signed descriptors and plural indexes.
  \item \textbf{Runtime and Provider Plane.} OpenFox, webTOS, third-party agent runtimes,
        model and tool services, storage, search and owner-operated edge terminals.
  \item \textbf{Application and Society Plane.} FreeCity, communication, software work,
        commerce and other intent-native applications and institutions.
\end{enumerate}

Only TOS Core and several reusable actors are substantially implemented as common foundations.
Profile-specific messaging and mailbox primitives exist, while the generic operation
registry, common request and receipt envelope, provider conformance suite and several economic
profiles remain partial or planned. This distinction is deliberate. Validators do not execute
application workloads, and adding a model, messenger feature, storage backend or commercial
profile must not require a consensus upgrade.

TOS does not seek to replace HTTP, MCP, A2A, OpenAPI, ARD or application-specific protocols.
Those systems can carry, describe or discover an interface. TOS adds portable identity,
authority and commitment around the action: who may propose, what was accepted, who may call,
within what limits, at what maximum cost, with what evidence and through what settlement path.

\subsection{Strategic Positioning}

The stable top-level positioning is:

\begin{quote}
\textbf{TOS Network is the foundational infrastructure---The Open System---for the Agentic
Internet.}
\end{quote}

The protocol-level expression is that TOS enables autonomous agents to discover,
communicate, coordinate, transact, execute and settle across organizational and platform
boundaries. The application-level expression is that an ordinary Internet action can become
typed, authorized, metered, receipted and settleable.

A useful architectural shorthand is that TOS provides the \textbf{economic syscall layer for
the Agentic Internet}. This does not mean a TOS validator executes every external action as a
kernel instruction. It means an agent can treat an independently provided service as a
well-defined operation with portable semantics, bounded authority, resource accounting and a
verifiable outcome.

HTTP standardized transport. DNS standardized naming. OpenAPI, MCP and A2A describe different
forms of interface and interaction. ARD makes resources discoverable. TOS aims to standardize
the economically accountable action that follows discovery.

The unit of coordination is therefore an accepted, policy-compliant Agreement that may be
fulfilled by one or more Agentic Service Operations. It is not an hour of unidentified
hardware, an unbounded API call or an unauthenticated message. A consumer and provider agree
on an exact operation revision, input and output contract, deadline, region, data-egress
policy, resource limit, maximum charge, evidence level and recourse. The provider or target
agent retains the right to apply local admission and safety policy, executes within the
accepted bounds, returns a signed receipt and settles according to the Agreement.

\subsection{Why Now}

Four curves are converging.

\begin{enumerate}
  \item \textbf{Agents can act.} Tool-using models and durable agent runtimes can now plan,
        call APIs and complete multi-step workflows with decreasing human supervision.
  \item \textbf{Interfaces are becoming machine-readable.} OpenAPI, MCP, A2A and emerging
        discovery formats reduce integration cost, but they do not provide one portable
        economic and authority contract.
  \item \textbf{Machine-native payment is becoming practical.} Programmable wallets,
        stablecoin settlement and payment-bearing HTTP patterns demonstrate that software can
        buy a service without a human checkout flow. Payment proves transfer; it does not by
        itself prove delegated authority, bounded consumption, correct execution or recourse.
  \item \textbf{Verification and abuse become the scarce resources.} As the marginal cost of
        generating requests and outputs falls, spam, unauthorized spending, unverifiable work
        and autonomous retry loops become more expensive for recipients and operators.
\end{enumerate}

Public a16z crypto analysis describes the same market transition in adjacent terms: agents
are becoming economic actors, while portable identity, permissions, payment and trust remain
fragmented.
\footnote{Christian Crowley et al., ``The missing infrastructure for AI agents: 5 ways
blockchains can help,'' a16z crypto, April 16, 2026:
\url{https://a16zcrypto.com/posts/article/5-ways-blockchains-help-ai-agents}. This reference
provides market context only and does not imply endorsement of TOS Network.}
The TOS thesis is that the missing product is not another model or payment button. It is the
enforceable contract around an autonomous action.

\subsection{Beachhead, Expansion and Network Flywheel}

The architecture is horizontal, but the go-to-market sequence must be narrow. The first
external acceptance target is \textbf{machine-checkable software work}: independent agents
form an Agreement over an exact repository state, bounded task, tests, deliverables, evidence
and recourse, then execute and verify the outcome without sharing one runtime. This profile
exercises the full Intent-to-Settlement lifecycle and produces evidence that independent
implementations can inspect.

The first commercial service wedge is \textbf{paid model and tool invocation}: an agent
discovers an independent provider, receives a live quote, delegates a bounded budget, invokes
a typed operation, verifies a signed usage receipt and settles without creating a bilateral
account integration. This wedge has an immediate buyer, a measurable unit of work and a clear
comparison against API keys, monthly invoices and payment-only alternatives.

Messenger and Mailbox form the second reference family. They are strategically useful because
they stress-test identity, recipient policy, delegation, retention, delivery evidence,
refundable attention bonds and machine-scale spam resistance. They should validate the common
operation model after the compute wedge rather than divide initial go-to-market focus.

The reference applications built by TOS must be the protocol's first demanding customers.
They should use the same public envelopes, SDKs and receipts offered to third parties; an
internal privileged path would hide integration cost and weaken the evidence for adoption.
External design partners should include at least one agent builder and multiple unrelated
service providers whose requirements can shape the first frozen profile.

If these beachheads work across independent implementations, the intended flywheel is:

\begin{center}
\texttt{open operation definitions -> competing providers -> portable receipts}\\
$\longrightarrow$ \texttt{better matching, trust, price and availability -> more demand}\\
$\longrightarrow$ \texttt{more providers and operation profiles}.
\end{center}

This follows a broad--narrow--broad market-design sequence: define the horizontal primitive,
win a high-value transaction category, then expand through adjacent profiles. Open protocols
compound through third-party embedding, so integrations and conformance matter more than a
closed catalog.
\footnote{Scott Duke Kominers, ``3 takeaways on market design for web3 builders,'' a16z
crypto, September 19, 2024:
\url{https://a16zcrypto.com/posts/article/3-takeaways-on-market-design-for-web3-builders}.}

\subsection{Defensibility and Value Capture}

The defensible asset is not token issuance, a proprietary model or a closed list of services.
It is the interoperable graph of operation identities, exact revisions, provider descriptors,
capabilities, quotes, receipts, disputes, settlement history and measured reliability. Every
additional conforming provider increases choice; every additional verified receipt improves
admission and matching; every additional integration makes the operation vocabulary more
useful to agents elsewhere.

Openness is therefore a distribution strategy, not the absence of a moat. Durable advantage
must come from becoming embedded in agent runtimes and provider SDKs, producing the safest and
lowest-friction completion path, and accumulating credible cross-provider performance data.
If providers can multi-home with no loss of history, trust or liquidity, the network effect
will be weak. If agents can carry identity, policy and receipts while providers compete on the
same definitions, TOS can become neutral coordination infrastructure rather than one more
marketplace application.

TOS Core captures protocol demand through transaction and message fees, persistent actor
state, escrow and dispute execution, registry and role bonds, receipt or channel commitments,
and validator security. Service-provider revenue remains distinct. The network should permit
applications to quote familiar units or use sponsored and batched payment policies where
appropriate; TOS must earn default settlement through predictable cost, finality, developer
experience and credible neutrality rather than by making every payload or price TOS-only.

\subsection{Why a Dedicated TOS Settlement Layer}

A new blockchain is not justified merely because agents can transact. General-purpose chains,
stablecoins and payment-bearing HTTP can settle many calls. TOS Core is justified only if its
actor-native asynchronous execution, value-bearing messages, persistent agent policy,
parallel Service Actors and sharded settlement make the full authorization-to-receipt flow
safer, cheaper or easier to operate across independent parties.

That claim is falsifiable. Reference profiles must publish end-to-end latency, failure and
recovery behavior, settlement cost as a share of operation value, channel or batch efficiency,
state growth, validator decentralization and equivalent deployment comparisons. If a
general-purpose chain plus the open TOS envelopes performs as well, the operation protocol may
still be useful, but the investment case for TOS Core and the native asset is materially
weaker.

\subsection{Falsifiable Adoption Gates}

Before describing this system as product-market fit or a network effect, public reporting
should show at least:

\begin{itemize}
  \item weekly and monthly cohorts of unrelated active payers, providers and agent builders;
  \item quote-to-paid-completion conversion, failed-match rate and time to successful match;
  \item payer and provider retention without token rewards or protocol-funded demand;
  \item organic settled value and its share of all settled value;
  \item service revenue, protocol fees and subsidies as separately reported economic flows;
  \item operation completion, timeout, cancellation and dispute-overturn rates;
  \item resource-estimate error between quoted limits and actual metered use;
  \item receipt verification and cross-provider conformance rates;
  \item the share of operations and providers reached through third-party integrations;
  \item spam rejection, bond refund and false-positive rates by communication profile;
  \item concentration of settled value and AIPoW score after common-control aggregation; and
  \item issuance divided by organic settled value when AIPoW is active.
\end{itemize}

Subsidized circular volume, address count, message count and declared capacity are not
substitutes for these measures. Section~\ref{sec:aipow} defines the demand and reward
quantities used by AIPoW reporting.

\subsection{Scale of the Opportunity}

If autonomous software continues to multiply, machine-to-machine operations become an
economic layer distinct from human-facing digital commerce. The addressable activity is not
limited to model inference. It includes communication, storage, event delivery, tools,
booking, procurement, verification, human work and physical services. TOS Network's design
goal is to provide the identity, authorization, metering, receipt and settlement
infrastructure this layer requires.

The venture-scale outcome is therefore not ``another AI marketplace.'' It is an open
transaction fabric used by many agent products and service networks, with TOS Core serving as
a preferred neutral settlement domain. The opposite outcome is also explicit: if agents stay
inside a few vertically integrated platforms, if payment-only protocols absorb the authority
and receipt layer, or if buyers do not value portable recourse, the addressable independent
network will be smaller. This is a structural thesis, not a claim about currently realized
transaction volume, protocol revenue or token value.

\section{From Intent to Settlement}
\label{sec:operation-model}

TOS treats an autonomous workflow as a sequence of distinct commitments. Collapsing these
states into one ``agent request'' makes it difficult to determine when authority began, what
was actually accepted, which evidence is relevant or whether payment is final.

\begin{center}
\texttt{Identity -> Discovery -> Intent -> Negotiation -> Agreement}\\
\texttt{-> Bounded Execution -> Evidence/Receipt -> Settlement}
\end{center}

\subsection{Identity and Discovery}

Identity identifies an agent, owner, provider, verifier or other principal and binds the keys,
roles and delegation policy relevant to the workflow. An agent identity must not imply
unrestricted owner authority. Key rotation, revocation, scope, expiry and delegation depth are
part of the authority model.

Discovery returns possible counterparties, capabilities, endpoints and operation profiles.
It is informative rather than authoritative: appearing in a catalog, index, DHT record or
search result does not prove current capacity, price, permission, trustworthiness or safety.
Clients re-verify signed descriptors, capabilities, quotes and settlement destinations before
committing funds or work.

\subsection{Intent, Negotiation and Agreement}

An \textbf{Intent} is a revocable expression of a desired outcome and constraints. It may name
a task, deadline, budget ceiling, privacy policy, preferred evidence or acceptable providers,
but it is not itself an authorization to spend or execute and does not create an obligation
for another party.

\textbf{Negotiation} converts an Intent into concrete proposals. Parties or their authorized
agents can refine scope, price, inputs, outputs, resource ceilings, deadlines, cancellation,
evidence, dispute forum and settlement method. A model may interpret natural-language goals
and propose terms; deterministic policy must decide whether those terms fit the principal's
authority and budget.

An \textbf{Agreement} is the exact, authenticated set of terms accepted by the relevant
parties. It binds identities, roles, selected operation revisions, capability references,
payload or schema commitments, budgets, evidence policy, recourse and expiry. An Agreement
may authorize one operation, a sequence, a conditional workflow or a reusable allowance. No
implementation should infer acceptance merely from discovery, conversation, possession of an
endpoint or an unsigned model output.

\subsection{Bounded Execution, Evidence and Settlement}

Agreement does not require validators to execute the real-world workload. Providers perform
model inference, communication, storage, software work, database access and physical action
off-chain under local admission, security and safety policy. Agentic Service Operations are
the typed execution primitives used to carry out the accepted terms. Every externality---cost,
compute, storage, delivery attempts, attention, retries or physical reach---is bounded by
capability, quota, budget, bond or another enforceable policy.

Evidence is evaluated only under the policy named by the Agreement. A Receipt records a
signed statement about status, metered use, result commitments, evidence references and
settlement state; it proves no more than its schema, signer and evidence policy authorize it
to prove. A stored receipt is not a delivery receipt, an execution receipt is not necessarily
a correctness proof, and a payment receipt is not proof that the underlying result was valid.

Settlement applies the accepted payment, refund, bond, dispute and finality rules. It may use
direct TOS transfers, escrow, prepaid balances, channels, batches, sponsorship or a zero-price
allowance. The invariant is accountable completion, not a mandatory per-call fee.

\subsection{Why an API Call Is Not an Agent Contract}

A conventional API call primarily specifies endpoint syntax and transport behavior. The
caller may still depend on provider-specific accounts, OAuth scopes, billing, rate limits,
error meanings, retry policy and evidence. Two APIs that both ``send a message'' or ``invoke a
model'' can be economically and operationally incompatible.

An Agentic Service Operation adds the information required for autonomous action across
organizational boundaries. It binds the requested behavior to identity, authority, resource
limits, economic policy and a verifiable result.

\begin{center}
\texttt{API call + capability + budget + metering + receipt + settlement}\\
$=$\\
\texttt{Agentic Service Operation}
\end{center}

The operation abstraction is intentionally transport-neutral. It may be carried through
RLDP/TOS Sites, HTTPS, QUIC, WebSocket, a message queue, a local IPC channel or another agreed
binding. Transport success does not imply operation authorization or settlement success.

\subsection{Agentic Service Opcode}

An \textbf{Agentic Service Opcode} is the architectural name for a versioned service-layer
action such as:

\begin{center}
\opid{messenger.send@1}\quad
\opid{mailbox.deposit@1}\quad
\opid{model.infer@2}\quad
\opid{storage.put@1}
\end{center}

The normative protocol field is \texttt{operation\_id}. ``Opcode'' is useful because the
operation behaves like a syscall from the agent's perspective: it has defined input, authority
requirements, cost and return semantics. It is not a claim that the action is a native TVM
instruction or that validators perform the external workload.

Multiple independent providers may implement the same operation revision. An agent can
discover candidates, compare quotes and evidence policy, authorize one within a bounded
budget, receive a standardized receipt and switch providers without learning a completely
different economic protocol.

\subsection{Three Opcode Layers}

This paper distinguishes three different meanings of opcode.

\begin{longtable}{L{0.22\textwidth}L{0.23\textwidth}L{0.21\textwidth}L{0.23\textwidth}}
\toprule
Layer & Examples & Execution boundary & Pricing boundary \\
\midrule
\endhead
TVM instruction opcode &
\texttt{SENDMSG}, \texttt{SETCODE} &
Deterministic validator execution &
Gas, cells, message forwarding and persistent storage \\
Contract message opcode &
\texttt{AGENT\_TASK\_SEND} and contract ABI operations &
Receiving actor interprets the message body &
Contract execution, attached value and message fees \\
Agentic Service Opcode &
\texttt{messenger.send@1}, \texttt{model.infer@2} &
External service, relay, terminal or another agent &
Profile-specific usage, quote, allowance, bond and settlement \\
\bottomrule
\end{longtable}

The layers may cooperate in one workflow. An off-chain
\opid{mailbox.deposit@1} operation can use an Agent Account for authorization, a Service Actor
or escrow for payment, and ordinary TVM instructions for the on-chain state transition.
They remain different abstractions and must not share ambiguous documentation or numeric
registries.

\subsection{Open Identity, Namespace and Versioning}

A global centrally assigned 32-bit number is not sufficient for an open Internet-scale
operation namespace. The full operation identity should be globally unambiguous, versioned
and bound into signatures. Examples include:

\begin{center}
\texttt{urn:tos:op:messenger.send:1}\\
\texttt{urn:tos:op:mailbox.deposit:1}\\
\texttt{urn:tos:op:example.com:weather.forecast:3}
\end{center}

The exact syntax is specification work, but the following rules are required:

\begin{itemize}
  \item the full operation identifier and revision are authenticated;
  \item incompatible semantics require a new immutable revision;
  \item schemas, metering dimensions, receipts and error meanings are content-addressed;
  \item domain-anchored or otherwise verifiable namespaces permit open extension;
  \item no single registry is mandatory for all operation publication; and
  \item compact numeric wire codes may be negotiated only inside an authenticated session
        after the full identifier has been agreed.
\end{itemize}

An operation definition is not a provider advertisement. The operation says what behavior and
semantics mean. A provider descriptor says that a particular identity implements that
operation under a particular endpoint, quote policy and evidence level.

\subsection{Formal Operation Definition}

A versioned operation definition can be represented as

\[
  \mathcal{O} =
  (id, v, S_{\mathrm{in}}, S_{\mathrm{out}}, K, M, R, E),
\]

where:

\begin{itemize}
  \item $id$ is the operation identifier and $v$ its immutable revision;
  \item $S_{\mathrm{in}}$ and $S_{\mathrm{out}}$ are canonical input and output schemas;
  \item $K$ defines capability, delegation and target-policy rules;
  \item $M$ defines resource dimensions and metering;
  \item $R$ defines receipt, evidence and settlement semantics; and
  \item $E$ defines errors, cancellation, retries and terminal states.
\end{itemize}

A request can be modeled as

\[
  Q =
  (op, caller, target, capability, nonce, expiry,
   \mathbf{u}_{\max}, C_{\max}, h(payload), correlation),
\]

where $\mathbf{u}_{\max}$ is a vector of resource ceilings and $C_{\max}$ is the maximum
authorized charge. The payload may remain encrypted and off-chain; the authenticated
envelope binds its digest.

A provider quote can be represented as

\[
  \mathcal{Q} =
  (provider, op, v, \mathbf{p}, \mathbf{u}_{\max},
   C_{\max}, t_{\mathrm{expiry}}, policy),
\]

where $\mathbf{p}$ prices the operation's resource dimensions. A receipt can be represented as

\[
  R =
  (status, \mathbf{u}_{\mathrm{actual}}, C_{\mathrm{actual}},
   h(result), evidence, settlement, \sigma_{\mathrm{provider}}).
\]

For accepted operations,

\[
  C_{\mathrm{actual}} \leq C_{\max}.
\]

A profile may also define minimum charges, reservation charges, cancellation fees, refunds,
sponsorship or a refundable bond. The accepted quote and operation definition determine the
meaning; a receipt cannot invent new dimensions after execution.

\subsection{Operation Execution Lifecycle}

A common lifecycle allows different profiles to share tooling and reasoning:

\begin{center}
\texttt{QUOTE}\\
$\downarrow$\\
\texttt{AUTHORIZE}\\
$\downarrow$\\
\texttt{RESERVE BUDGET}\\
$\downarrow$\\
\texttt{EXECUTE}\\
$\downarrow$\\
\texttt{RECEIPT}\\
$\downarrow$\\
\texttt{VERIFY}\\
$\downarrow$\\
\texttt{SETTLE}
\end{center}

This is the execution sub-lifecycle of an accepted Agreement, not a substitute for Intent and
negotiation. Every step is asynchronous. A quote may expire. Authorization may fail after discovery.
Admission may reject an otherwise affordable call. Execution may return a partial result.
Receipt verification may trigger dispute instead of settlement. These are explicit states,
not hidden exceptions.

Long-running operations require stable operation identity, idempotency keys and restart
recovery. Retrying a request must not silently duplicate a payment, message, reservation,
physical action or task.

\subsection{Metering and Charging Are Distinct}

Every meaningful operation must have bounded and attributable resource use, but not every
operation must charge a non-zero amount.

A valid policy may use:

\begin{itemize}
  \item direct payment in TOS;
  \item prepaid credit or subscription allowance;
  \item organization, recipient or application sponsorship;
  \item a trusted-party free quota;
  \item congestion-dependent pricing;
  \item a refundable anti-spam or performance bond;
  \item batched settlement or a payment channel; or
  \item deferred aggregate settlement.
\end{itemize}

The invariant is not ``every action pays.'' The invariant is:

\begin{quote}
\textbf{No caller may impose unbounded external cost without identity, authorization, quota,
budget or bond.}
\end{quote}

Even a zero-price operation remains metered. Its provider can account for bytes, time,
delivery attempts, compute, storage or recipient attention; enforce limits; produce a receipt;
and attribute abuse to an identity and capability.

\subsection{Representative Operation Profiles}

\begin{longtable}{L{0.22\textwidth}L{0.31\textwidth}L{0.37\textwidth}}
\toprule
Operation & Principal metering dimensions & Typical control and settlement policy \\
\midrule
\endhead
\opid{messenger.send@1} &
Ciphertext bytes, recipients, delivery attempts, priority &
Recipient capability, contact quota, introduction bond, relay quote \\
\opid{mailbox.deposit@1} &
Bytes, retention duration, replica count &
Prepaid storage allowance, capacity ceiling, signed stored receipt \\
\opid{mailbox.read@1} &
Requests and outbound bytes &
Read capability, device binding, rate limit \\
\opid{mail.send@1} &
Body and attachment bytes, delivery class &
Sender reputation, domain policy, refundable attention bond \\
\opid{model.infer@1} &
Input and output tokens, accelerator time, latency class &
Maximum budget, provider quote, usage and evidence receipt \\
\opid{tool.invoke@1} &
Calls, duration, external provider cost &
Capability, deadline, escrow, normalized result and errors \\
\opid{storage.put@1} &
Bytes, duration, redundancy and egress &
Lease, prepayment, deletion and durability policy \\
\opid{event.publish@1} &
Events, subscribers, fan-out and retention &
Publisher quota, congestion price and bounded replay \\
\opid{booking.reserve@1} &
Resource quantity and lock duration &
Refundable deposit, cancellation window and confirmation receipt \\
\opid{task.post@1} &
Budget, deadline and verification tier &
Escrow, claimant bond, result receipt and dispute policy \\
\bottomrule
\end{longtable}

These identifiers are illustrative until frozen by their respective specifications. Their
purpose is to show that communication, compute, storage and commerce can share one economic
operation model without sharing one workload implementation.

\section{Design Principles}

\subsection{Typed Operations, Not Bare Endpoints}

Agents request a versioned operation and declared outcome. A URL, host, model name, device
class or raw hardware claim may help discovery and scheduling, but it is not the operation
contract. Providers expose bounded service behavior rather than an unrestricted shell,
container socket, accelerator or actuator.

\subsection{Actor First}

Every account is an actor. Agent Accounts, Task Escrows, capability entries, service actors
and verifier actors are ordinary native accounts with isolated state. There is no privileged
global application process.

\subsection{Asynchronous by Default}

Cross-actor and cross-provider workflows use messages, callbacks, deadlines, receipts and
retries. A sender cannot assume synchronous execution or multi-step atomicity. Contracts and
service profiles expose explicit intermediate states and deterministic failure behavior.

\subsection{Authority Must Be Explicit}

Owner, controller, creator, assigned agent, recipient, verifier, service, relay and fee-payer
roles are distinct. Every state transition, resource reservation and transfer identifies its
authority. Discovery, possession of an endpoint and payment capacity do not grant execution
authority.

\subsection{Models Interpret; Deterministic Policy Authorizes}

A model may translate natural language into an Intent, propose an execution plan or evaluate
non-deterministic evidence. It must not silently expand capability, budget, target, delegation
depth or physical reach. The accepted Agreement and deterministic policy checks remain the
authority boundary. Ambiguity, unavailable evidence or policy conflict produces refusal,
renegotiation or human escalation rather than inferred permission.

\subsection{Meter the Externality}

The initiator should fund or otherwise hold authorization for the resource and attention cost
it imposes. Bytes, storage duration, delivery attempts, compute, tool expense, reservation
time and retries are profile-specific resources. A free allowance is still finite and
attributable.

\subsection{Verification Without Payload Bloat}

The chain is authoritative for funds, permissions, deadlines and settlement where shared
trust is required. Large messages, model prompts, outputs, datasets, sensor streams and
private credentials remain off-chain. Contracts store hashes, signatures, attestations and
evidence references that bind external work to accepted operations.

\subsection{Local Control, Privacy and Safety}

Providers retain custody of hardware, data, models and operational policy. Encrypted
communication payloads and private execution context remain outside public consensus. A
physical terminal continues approved local work while disconnected and keeps blockchain
consensus outside hard real-time and actuator-control loops. Participation in TOS does not
transfer ownership or control of a provider's infrastructure; the provider may reject work
that violates current admission, legal, security, capacity or safety policy even when a quote
or earlier Agreement exists, subject to the agreed cancellation, refund and recourse rules.

\subsection{Bounded Resources and Recoverable Failure}

Connections, queues, retries, subscriptions, message retention, model caches, RAM,
accelerator memory, durable journals and watchers have explicit limits and cleanup paths.
Failure, cancellation, restart, disconnect and payment rejection are normal protocol states.
No unauthenticated, unbudgeted, unmetered or policy-disallowed caller may cause unbounded
memory, disk, compute, delivery, retry or durable-state growth.

\subsection{Native Execution Focus}

TOS Core has one active application execution domain: the native TVM basechain. External
service workloads run in separate agents, relays, terminals and provider repositories, not
inside \texttt{validator-engine}. New operation profiles must not require validators to embed
application-specific runtimes.

\subsection{Open Discovery and Provider Competition}

TOS uses open discovery mechanisms and does not require one mandatory catalog. A service
operation may be implemented by multiple providers and carried over multiple transports.
Critical clients re-verify identity, capability, quote, endpoint and settlement state rather
than treating a search index as authority.

\section{System Architecture}

\subsection{Five-Plane Infrastructure Architecture}

TOS Network is designed as five interoperable planes with independently released components.
The separation limits consensus risk, permits protocols and applications to evolve at their
appropriate speed and keeps private workloads outside validator execution. A plane is a
responsibility boundary, not a claim that every component is implemented or operated by one
organization.

\begin{longtable}{L{0.18\textwidth}L{0.35\textwidth}L{0.36\textwidth}}
\toprule
Plane & Responsibility & Delivery boundary \\
\midrule
\endhead
Trust and Settlement &
TOS Core consensus, TVM, actor accounts, value-bearing messages, contracts, assets, escrow,
disputes, commitments and finality &
Implemented foundations in the \texttt{tos} repository; production readiness remains subject
to the stated release gates \\
Identity and Operation &
Agent and provider identity, ownership, capabilities, Intent, Agreement, versioned operations,
authorization, replay protection, resource envelopes, metering, evidence and Receipts &
Implemented reusable actors plus partial or planned open specifications, SDKs and conformance \\
Propagation and Discovery &
Direct communication, ADNL/DHT/RLDP, TOS DNS, TOS Sites, Messenger, Mailbox, ARD publication,
gateways, signed descriptors, indexes and federated registries &
Core network primitives implemented; common propagation, discovery and interoperability
products are partial or planned \\
Runtime and Provider &
OpenFox, webTOS, third-party agent runtimes, model and tool services, storage, search,
relays and owner-operated physical terminals &
Independently released runtimes, adapters and operator infrastructure; providers retain local
admission, pricing, data and safety authority \\
Application and Society &
FreeCity, communication, software work, commerce, booking, human services, communities and
other intent-native applications and institutions &
Application-specific state machines and interfaces built above the public infrastructure \\
\bottomrule
\end{longtable}

Existing Agent, Service, Escrow, Registry, Dispute and Attestation actors remain reusable TOS
Core foundations. Product-specific daemons, message payloads, model runtimes, catalogs and
continuous device telemetry do not belong in the validator repository. A Gateway, Carrier or
Indexer may improve reachability and usability, but none is protocol authority for identity,
capability, accepted Agreement, funds or final settlement.

\subsection{The Economic Syscall Boundary}

From an agent's perspective, an Agentic Service Operation resembles a syscall: it has a stable
name, typed arguments, authority requirements, resource limits, a defined return status and an
accountable cost. Unlike a host operating-system syscall, it may cross administrative,
network and legal boundaries and may be implemented by competing providers.

The shared TOS boundary covers:

\begin{itemize}
  \item persistent identities and delegated authority;
  \item operation and schema commitments;
  \item quotes, budgets, escrow and payment authorization;
  \item receipt and evidence commitments;
  \item disputes, refunds and final settlement; and
  \item reputation or bond state where a profile defines it.
\end{itemize}

The provider boundary covers:

\begin{itemize}
  \item admission under local policy;
  \item private payload handling and decryption;
  \item actual communication, storage, inference, tool use or physical execution;
  \item detailed metering and operational logs;
  \item result generation; and
  \item signing the agreed receipt and evidence envelope.
\end{itemize}

\subsection{On-Chain, Off-Chain and Derived State}

\begin{longtable}{L{0.25\textwidth}L{0.30\textwidth}L{0.34\textwidth}}
\toprule
State class & Representative content & Authority \\
\midrule
\endhead
On-chain shared state &
Balances, Agent Account policy, capability and descriptor commitments, escrow, deadlines,
receipt commitments, dispute state and settlement &
Consensus-authoritative where the applicable contract or protocol rule says so \\
Off-chain private or operational state &
Encrypted messages, prompts, model outputs, files, sensor streams, provider logs, queue state
and local safety policy &
Authoritative only under the accepted operation and evidence policy \\
Derived index and analytics &
Search indexes, rankings, reconstructed timelines, provider statistics and dashboards &
Non-authoritative; critical clients re-verify signed or on-chain state \\
\bottomrule
\end{longtable}

Not every operation needs an on-chain transaction. Trusted allowances, payment channels,
prepaid balances and batched settlement may carry high-frequency actions. The chain is used
where independently operated parties need shared authority, dispute or final economic state.

\subsection{Masterchain and Basechain}

The masterchain anchors validator sets, consensus configuration, workchain descriptors and
cross-shard finality. The basechain executes native TVM accounts and contracts. The canonical
genesis registers the masterchain with identifier $-1$ and the native basechain with identifier
$0$.

The masterchain does not execute a model, deliver a private message or perform an external
service workload. It establishes the shared state required for agents and service actors to
trust identity, authority, account state and settlement outcomes.

\subsection{Accounts as Actors}

An account transition can be represented as

\[
  (S_{t+1}, M_{out}) = F(S_t, M_{in}, C_t),
\]

where $S_t$ is the account state, $M_{in}$ is one inbound message, $C_t$ is the consensus
execution context and $M_{out}$ is a finite set of outbound messages. The transition changes
only the receiving account. Effects on another account occur later when an emitted message is
delivered there.

This isolation maps directly to agentic systems: planners, workers, tools, tasks, relays and
verifiers can progress independently while retaining a shared, auditable settlement layer.

\subsection{Cells and TVM}

TOS stores account code and persistent state in immutable cells. Cells form directed acyclic
graphs and are addressed by cryptographic hashes. TVM executes contract code deterministically,
charges gas and emits messages. Compact cell-native state is well suited to balances, policy,
status values and content commitments.

\subsection{Asynchronous Messages}

Messages carry a destination, optional source, value and body. Internal messages have an
unforgeable source account. External messages may carry signatures and replay-protection data.
Value transfer is therefore part of the same mechanism as task acceptance, service calls and
settlement.

Contract messages should include a contract message opcode, a correlation value such as
\texttt{query\_id}, clear value semantics and enough domain binding to prevent replay against
another account, task or network. When the message authorizes an external Agentic Service
Operation, it also binds the full \texttt{operation\_id}, revision, capability, budget and
payload or quote commitment. The contract message opcode and service operation identifier are
different layers.

\subsection{Sharding}

TOS follows a bottom-up sharding model. Logically, each account has an independent history;
physically, many account histories are grouped into shardchains. Shards may split and merge as
load changes. Cross-shard messages preserve the actor abstraction even when communicating
accounts reside in different shards.

This model permits parallel execution because transactions affecting different destination
accounts do not mutate shared application state. High-volume agent economies can distribute
tasks and services across many accounts instead of concentrating all workflow state in one
contract.

\subsection{Consensus and Networking}

Validators agree on masterchain and shardchain blocks and enforce protocol configuration.
The networking stack provides authenticated datagram transport, distributed discovery,
reliable large-message delivery and overlay communication. ADNL identities and TOS DNS can
locate services; RLDP and TOS Sites can carry application traffic. These primitives do not by themselves provide the planned operation definition, quote,
receipt, Messenger relay, Mailbox, model service, storage service or terminal product.
\texttt{validator-engine} remains the canonical node process; \texttt{tosctl} provides an
operator path for configuration and existing agent-oriented workflows.

\section{Trust and Settlement Actors}

\subsection{Agent Wallet}

An Agent Wallet is the custody and local-profile layer for an autonomous agent. The local
profile records the underlying native wallet, owner key reference, controller key reference,
spending policy, capabilities, optional metadata commitments and optional runtime binding.

Creating a profile generates separate owner and controller keys in the configured secrets
vault. Funding transfers TOS to the derived native wallet address. Activation deploys the
underlying wallet contract after the address has sufficient balance.

The owner retains full custody authority. Manual owner transfers and emergency operations are
operator actions. Automated agent spending must use the policy-enforced Agent Account path.
The owner key must never be exported to an off-chain agent runtime.

\subsection{Agent Account}

An Agent Account is the on-chain authority boundary for automated actions. Its state includes:

\begin{itemize}
  \item owner address and controller public key;
  \item sequence number and expiry checks for replay protection;
  \item maximum value per controller action;
  \item cumulative daily limit and current daily spend;
  \item default task timeout;
  \item optional metadata and service-endpoint hashes.
\end{itemize}

Owner-authorized internal messages update policy or rotate the controller. Controller-signed
external messages may send bounded task actions only after signature, expiry, sequence,
per-action and UTC-day checks succeed. Local CLI checks improve operator feedback, but the
contract is the final authorization boundary.

\subsection{Task Actor}

A Task Escrow is an account representing one unit of work. It stores the creator, optional
assigned agent, optional verifier, budget, deadline, review period, lifecycle status, result
hash, evidence hash, settlement-policy hash, permission hash and dispute hash.

The implemented lifecycle includes:

\begin{center}
\texttt{Open -> Accepted -> Submitted -> Settled}
\end{center}

with terminal or review paths for rejection, cancellation, expiry and dispute resolution. An
unassigned open task may be claimed atomically; the first valid claim records the caller as the
agent. Result submission starts a review window. Authorized settlement releases a bounded
payout and returns any remainder according to contract rules.

\subsection{Capability Registry Actor}

A Capability Registry entry is one contract per registered agent or service. It is not a
chain-wide dictionary and does not by itself provide complete network enumeration. Its state
commits to:

\begin{itemize}
  \item owner and optional verifier;
  \item active or inactive status and registration time;
  \item task-category, pricing, metadata and verification-method hashes;
  \item a stakeable bond;
  \item an accumulating signed reputation score.
\end{itemize}

Owner-authorized operations update metadata, rotate the verifier, withdraw bond, deactivate
or reactivate the entry. Staking is permissionless. Reputation updates require verifier
authority. The registry exposes inspectable commitments while allowing detailed service
descriptions to remain off-chain.

\subsection{Service Actor}

A Service Actor represents a bounded implementation of a model, data, tool, relay, storage,
communication or other defined Agentic Service Operation. The implemented
contract supports open access or one authorized caller, price per call, a UTC-day rate limit,
metadata and proof-scheme hashes, request and response commitments, accumulated revenue and an
active or inactive lifecycle. A call must attach at least the configured price; accepted value
is credited to contract revenue. The owner commits response hashes, updates policy, withdraws
revenue and controls activation. The registry says what an operation provider claims to implement; the Service Actor governs
calls and payments where that profile uses an on-chain actor. It does not expose a host, raw accelerator or
arbitrary execution environment.

\subsection{Verifier and Attestation Actors}

A Verifier Actor evaluates result or evidence commitments according to a declared policy. A
verifier may accept, reject, score or dispute work, but its authority must be explicitly stored
by the task or registry actor. The implemented Proof Attestation foundation can record compact
claims and evidence references. Neither a payment nor a hash alone proves semantic
correctness, confidentiality or the physical hardware used. TOS does not hard-code one model,
oracle or proof backend.


\subsection{Operation Definition and Provider Descriptor}

An operation definition and a provider descriptor are distinct.

The operation definition commits the portable semantics: identifier, revision, schemas,
capability rules, metering dimensions, receipt requirements and error meanings. A provider
descriptor binds a particular identity to supported operation revisions, endpoints, pricing
or quote method, evidence level, region and current authorization.

Operation definitions may be published through domain-anchored namespaces, signed registries
or content-addressed repositories. Provider descriptors may refer to on-chain Capability
Registry and Service Actor commitments. Neither a search result nor an unsigned descriptor
grants spending or execution authority.

A generic on-chain Operation Registry is planned rather than implied by the current
Capability Registry. The open namespace must allow independent publication and mirrored
indexes without placing all Internet operation definitions in one chain-wide dictionary.
\section{Agent Authority and Delegation Model}

\subsection{Owner and Controller Separation}

The owner controls custody, recovery and high-authority policy changes. The controller is an
automation key with bounded authority. Compromise of a controller should not become
owner-equivalent compromise.

For a controller action with value $v$, per-action limit $L_a$, daily limit $L_d$ and already
spent amount $D$, the Agent Account requires

\[
  0 < v \leq L_a \quad\text{and}\quad D + v \leq L_d.
\]

It also verifies a valid signature, an expected sequence number and a non-expired timestamp.

\subsection{Operational Lifecycle}

The initial operator lifecycle is:

\begin{enumerate}
  \item create a local Agent Wallet profile and vault keys;
  \item fund and activate the underlying native wallet;
  \item build and inspect deterministic Agent Account state;
  \item deploy the Agent Account through an active funding wallet;
  \item bind an off-chain runtime to the profile;
  \item execute bounded task actions with the controller;
  \item use the owner for policy updates, controller rotation and explicit custody actions.
\end{enumerate}

\subsection{Policy and Permission Linkage}

Task actions can bind a locally tracked permission identifier to an on-chain permission hash.
Before signing, tooling checks the target contract, lifecycle state, assignment and permission
linkage. The Task Escrow independently validates sender authority and state transition rules.
This layered design catches operator mistakes without treating local configuration as chain
authority.

\section{Task Markets and Settlement}

\subsection{Creating Work}

A creator deploys a Task Escrow with a budget, deadline, settlement policy and optional agent
or verifier. The deployment value must cover the intended escrow and execution costs. Task
metadata itself may be stored off-chain; the contract retains compact commitments.

\subsection{Assignment and Claim}

A task can target a specific Agent Account or remain open. A targeted agent may accept or
reject according to contract state. An open task may be claimed atomically. Assignment is
authoritative only after it appears in chain state.

\subsection{Results and Evidence}

The agent submits a result hash and evidence hash. These values do not prove semantic quality
by themselves; they bind a later-retrieved artifact or attestation to the task. The review
period gives the creator or verifier time to settle or dispute the submission.

\subsection{Settlement and Disputes}

Settlement transfers no more than the available escrow. Cancellation and rejection refund
according to explicit lifecycle rules. Timeout is consensus-time based. A dispute records a
reason or evidence commitment and moves authority to the configured resolution path. Verifier
resolution is bounded by contract balance and role checks.


\section{Agentic Operation Protocol}
\label{sec:protocol}

\subsection{Protocol Scope}

The Agentic Operation Protocol is the common execution protocol within the Identity and
Operation Plane above TOS Core. Its purpose is not to define every business workflow. It
defines the portable contract through which an accepted Agreement authorizes one or more
operations so independent agents and providers can interoperate safely.

Independent implementations require canonical and versioned definitions for:

\begin{itemize}
  \item operation identifiers, revisions and schema hashes;
  \item signed provider descriptors and short-lived endpoint manifests;
  \item authentication, replay domains, nonces and bounded delegation;
  \item resource dimensions, limits, quotes and maximum charges;
  \item payment authorization, sponsorship, allowances, bonds and refunds;
  \item request, stream, result, receipt and evidence envelopes;
  \item durable operation identities, cancellation, retry and restart recovery;
  \item error classes, limits and compatibility rules; and
  \item batch, channel and final settlement references.
\end{itemize}

This protocol is planned common work, not functionality implied merely by the existence of a
Service Actor contract or one profile-specific implementation. The contract is a settlement
primitive; interoperable sessions and receipts still require schemas, SDKs, conformance
vectors and independent implementations.

\subsection{Request Envelope}

A canonical request envelope includes at least:

\begin{longtable}{L{0.27\textwidth}L{0.62\textwidth}}
\toprule
Field & Meaning \\
\midrule
\endhead
\texttt{operation\_id}, \texttt{revision} &
Exact operation semantics and immutable revision \\
\texttt{request\_id} &
Globally unique or replay-domain unique operation instance \\
\texttt{correlation\_id}, \texttt{causation\_id} &
Task or workflow lineage and immediate cause \\
\texttt{caller}, \texttt{target}, \texttt{fee\_payer} &
Identities participating in authorization and payment \\
\texttt{capability} &
Scope, delegation chain, quotas and target policy \\
\texttt{nonce}, \texttt{expiry} &
Replay and stale-request protection \\
\texttt{payload\_digest} &
Authenticated commitment to an inline or external, possibly encrypted payload \\
\texttt{resource\_limits} &
Profile-defined ceilings such as bytes, tokens, duration, attempts or retention \\
\texttt{max\_cost} &
Maximum authorized charge under the accepted asset and settlement policy \\
\texttt{quote\_ref} &
Accepted provider quote and expiry \\
\texttt{receipt\_policy} &
Required receipt, evidence and acknowledgement levels \\
\texttt{idempotency\_key} &
Stable duplicate-suppression identity across retry and restart \\
\bottomrule
\end{longtable}

The signed domain includes the network, operation identifier and revision, target identity,
capability, request identity, nonce, expiry, resource limits, maximum cost and payload digest.
A valid request for one operation, target, network or quote must not authorize another.

\subsection{Quotes and Admission}

Discovery returns candidates, not live capacity or permission. A provider returns a
short-lived quote for an exact operation revision and resource envelope. The quote identifies:

\begin{itemize}
  \item provider and payment destination;
  \item operation identifier and revision;
  \item priced resource dimensions and minimum charge;
  \item maximum admitted limits and current expiry;
  \item cancellation, reservation and refund policy;
  \item evidence and receipt level;
  \item region, privacy and data-egress policy; and
  \item any required bond, stake or prepaid allowance.
\end{itemize}

The provider remains authoritative for local admission. A valid quote does not force a
provider to violate updated safety, resource, legal or recipient policy. If admission fails,
the profile defines which reservation charges are retained or refunded.

\subsection{Capability and Delegation}

A capability grants a bounded right to request one or more operations. It includes:

\begin{itemize}
  \item issuer, subject and target scope;
  \item operation identifiers and permitted revisions;
  \item resource, frequency and cumulative spending limits;
  \item delegation depth and whether sub-delegation is permitted;
  \item valid time interval and revocation reference;
  \item optional recipient, organization, device or endpoint binding; and
  \item policy for sponsorship, bond and receipt visibility.
\end{itemize}

Capabilities follow least authority. A capability to deposit encrypted content into one
Mailbox does not authorize reading it, sending arbitrary messages, spending from the owner's
wallet or invoking another provider. A discovery record or public endpoint is not a
capability.

\subsection{Metering Schema}

Each operation revision defines a finite vector of resource dimensions. Examples include:

\[
  \mathbf{u} =
  (bytes,\ seconds,\ tokens_{\mathrm{in}}, tokens_{\mathrm{out}},
   attempts,\ replicas,\ recipients).
\]

Not every profile uses every dimension. Units, rounding, sampling, measurement authority and
maximum error must be specified. The quote fixes the price vector or a deterministic pricing
function before execution.

The actual charge may be modeled as

\[
  C_{\mathrm{actual}} =
  C_{\mathrm{admission}}
  + \sum_i p_i u_i
  + C_{\mathrm{external}}
  + C_{\mathrm{priority}}
  - C_{\mathrm{credit}},
\]

subject to the accepted maximum and refund policy. A provider cannot add an undocumented
resource dimension after execution. When measurement is not independently reproducible, the
receipt states the trust and evidence level.

\subsection{Receipt and Evidence Envelope}

A receipt is the provider's signed statement about one operation instance. A minimum receipt
includes:

\begin{itemize}
  \item request, operation and revision identifiers;
  \item caller, target, provider and fee-payer bindings;
  \item start, terminal and receipt timestamps;
  \item terminal status and normalized error class;
  \item actual resource vector and actual charge;
  \item result, payload or transcript digest;
  \item evidence level and references;
  \item refund, bond and settlement state; and
  \item provider signature and optional co-signatures.
\end{itemize}

Different receipts prove different facts. A \emph{stored receipt} says a Mailbox accepted
ciphertext under a retention policy. A \emph{delivery receipt} says a relay or device accepted
delivery. An \emph{acknowledgement receipt} says the target agent emitted an application-level
acknowledgement. A \emph{settlement receipt} says an economic state reached its specified
finality. None automatically proves the others.

A receipt proves that its signer made the stated claim. Semantic correctness still depends on
the accepted evidence and dispute policy.

\subsection{Normalized Status and Errors}

Profiles share a common outer status taxonomy even when their detailed errors differ:

\begin{center}
\texttt{ACCEPTED, RUNNING, PARTIAL, SUCCEEDED, REJECTED, CANCELLED,}\\
\texttt{EXPIRED, FAILED, DISPUTED, REFUNDED, SETTLED}
\end{center}

Stable error classes include invalid request, unauthorized, capability expired, quota
exceeded, quote expired, insufficient budget, admission rejected, unavailable, timeout,
cancelled, provider failure, evidence failure, settlement failure and incompatible revision.

Unknown critical extensions cause rejection. Unknown non-critical extensions may be ignored
only where the operation revision explicitly permits it.

\subsection{Idempotency, Retry and Recovery}

A provider persists enough operation state to distinguish:

\begin{itemize}
  \item an unseen request;
  \item an accepted but unfinished request;
  \item a completed request whose receipt must be replayed;
  \item a cancelled or expired request; and
  \item a conflicting request reusing an idempotency identity.
\end{itemize}

Retries are bounded, funded or quota-backed and use backoff. A reconnect must not send a
message, charge a payer, reserve a room, transfer an asset or actuate a device twice. Durable
journals have explicit size and age limits and compaction rules.

\subsection{Transport and Interface Bindings}

The operation protocol may bind to:

\begin{itemize}
  \item ADNL/RLDP and TOS Sites;
  \item HTTPS, HTTP/2 or HTTP/3;
  \item QUIC or WebSocket;
  \item MCP tools or resources;
  \item A2A agent endpoints;
  \item OpenAPI operations;
  \item local IPC; or
  \item profile-specific encrypted sessions.
\end{itemize}

A binding defines canonical serialization, authentication framing, streaming, cancellation
and error mapping. It does not change the operation identity or economic semantics.

\subsection{High-Frequency and Aggregate Settlement}

One on-chain transaction per message, token chunk or event is neither necessary nor desirable.
Profiles may use:

\begin{itemize}
  \item prepaid balances with provider-signed usage receipts;
  \item bounded bilateral or multi-party payment channels;
  \item organization sponsorship accounts;
  \item periodic receipt-root commitments;
  \item netted batch settlement; and
  \item refundable bonds settled only after an abuse or dispute outcome.
\end{itemize}

Final settlement remains bounded by the accepted authorization and auditable receipt set.
Channels and batches must define replay, sequence, closure, timeout and dispute behavior.
High frequency is not permission for unbounded provider credit or unbounded caller liability.

\section{Discovery, Naming and Interoperability}

\subsection{Division of Responsibility}

TOS is complementary to existing Internet and agent protocols.

\begin{longtable}{L{0.20\textwidth}L{0.33\textwidth}L{0.36\textwidth}}
\toprule
Layer or protocol & Primary role & What it does not grant by itself \\
\midrule
\endhead
DNS and TOS DNS &
Human-readable and machine-resolvable naming &
Capability, live capacity, spending or execution authority \\
ARD &
Open publication and search for agents, tools, APIs, workflows and catalogs &
Authentication, quote, payment, receipt or settlement \\
OpenAPI, MCP and A2A &
Interface or interaction binding for tools, APIs and agents &
Portable economic policy across independent providers \\
HTTP, QUIC, WebSocket and RLDP &
Transport and streaming &
Operation semantics, authority or final economic state \\
Agentic Operation Protocol &
Capability, resource budget, quote, metering, receipt and settlement contract &
The private workload implementation itself \\
TOS Core &
Shared identity, actor state, value, commitments, dispute and settlement &
Private payload delivery, model execution or physical control \\
\bottomrule
\end{longtable}

\subsection{Competitive Boundary}

TOS competes with combinations of existing layers, not with an empty market. The relevant
comparison is whether one integrated operation contract produces enough incremental safety,
portability and completion to justify another dependency.

\begin{longtable}{L{0.22\textwidth}L{0.29\textwidth}L{0.38\textwidth}}
\toprule
Alternative & Strength & Remaining TOS claim to prove \\
\midrule
\endhead
Payment-bearing HTTP and machine-payment protocols &
Low-friction payment attached to a request &
Portable delegation, cumulative budgets, multidimensional metering, evidence, normalized
failure, refunds and disputes across payment methods \\
MCP, A2A and OpenAPI ecosystems &
Broad developer adoption for tools, agents and interfaces &
One economic envelope that can bind to those interfaces without forking or replacing them \\
General-purpose L1s and L2s &
Existing liquidity, wallets, developers and composable contracts &
Lower total integration cost and better asynchronous agent-service execution, demonstrated
with comparable workloads rather than asserted from architecture \\
Vertical agent marketplaces &
Curated supply, managed trust, distribution and immediate user experience &
Portability across marketplaces and providers, with credible neutrality and no mandatory
single operator \\
Centralized API management and cloud billing &
Mature operations, support and familiar enterprise controls &
Cross-organization authority and recourse where no platform is accepted as the permanent
trusted intermediary \\
\bottomrule
\end{longtable}

The separation permits TOS to reuse open ecosystems rather than creating incompatible
replacements. A resource discovered through ARD can expose MCP, A2A, OpenAPI, HTTPS or RLDP
bindings while keeping the same TOS operation identity, capability and settlement semantics.

\subsection{TOS and Bittensor}

Bittensor is an important adjacent project, but the two systems begin from different primary
problems. Bittensor describes itself as a framework or language for open, decentralized
digital-commodity markets called subnets, organized under a shared token system and directed
ultimately toward machine intelligence. Its official documentation describes independent
subnets producing commodities such as compute, inference, storage and prediction: miners
produce the commodity, validators score miners, subnet creators define incentive mechanisms
and the chain pays participants in TAO according to contributed value.
\footnote{Bittensor, ``The Bittensor Paradigm,''
\url{https://www.bittensor.com/about}, and ``Bittensor Documentation,''
\url{https://www.bittensor.com/docs}, accessed August 29, 2026. Bittensor terminology and
functionality may evolve; this comparison uses those official descriptions and does not imply
endorsement by either project.}

TOS starts from the cross-agent coordination problem. It aims to provide identity,
authorization, communication, discovery, Intent, Agreement, bounded execution, evidence,
Receipts, dispute and settlement across independent agents and providers. Machine
intelligence and digital commodities can be supplied through this infrastructure, but they
are not the complete scope of the Agentic Internet.

\begin{longtable}{L{0.19\textwidth}L{0.34\textwidth}L{0.36\textwidth}}
\toprule
Dimension & Bittensor & TOS Network \\
\midrule
\endhead
Primary objective &
Open markets that incentivize production and evaluation of digital commodities, culminating
in machine intelligence &
Open foundational infrastructure through which autonomous agents form accountable
commitments and act across organizational boundaries \\
Primary coordination unit &
A subnet with its commodity, miners, validators and incentive mechanism &
An authenticated lifecycle from Intent and Agreement through one or more bounded Operations,
Evidence/Receipt and Settlement \\
Evaluation &
Subnet-specific validators score miners according to the incentive mechanism defined for the
commodity &
Agreement-specific evidence and verifier policies assess only declared claims; providers and
principals retain local admission and trust policy \\
Network organization &
Multiple independent commodity subnets connected through Bittensor's chain and TAO system &
Five interoperable planes spanning trust, identity, propagation, runtimes and applications,
with plural providers, indexes, gateways and transports \\
Blockchain role &
Registers and coordinates the network and pays participants in TAO in proportion to the value
recognized by its mechanisms &
Provides shared actor state, commitments, escrow, disputes and final settlement where
counterparties need a tamper-resistant source of truth \\
Application scope &
Digitally produced and evaluated commodities such as compute, inference, storage and
prediction &
Cross-agent communication and cooperation, including software work, commerce, booking,
human services and physical-edge activity as well as compute and inference \\
\bottomrule
\end{longtable}

The distinction is architectural, not a claim that the projects are mutually exclusive. A
Bittensor-backed inference, storage, prediction or other subnet service can publish a TOS
Capability and operation descriptor. A TOS agent can then discover that service, negotiate
terms, form an Agreement, invoke it through its supported binding, collect the agreed evidence
and Receipt, and settle under the chosen policy. In that composition, Bittensor can organize
and incentivize a commodity network while TOS supplies the cross-agent identity, authority,
commitment and accountability envelope around its use.

\subsection{Agentic Resource Discovery Compatibility}

TOS adopts Agentic Resource Discovery (ARD) as the public, protocol-neutral discovery
envelope for agents, APIs, tools, workflows and nested catalogs. The initial compatibility
target is ARD v0.9 Draft as assessed on July 31, 2026.\footnote{Agentic Resource Discovery
specification: \url{https://agenticresourcediscovery.org/spec/}.} Because that document is a
proposal rather than a frozen standard, every implementation records the exact version it
accepts and does not silently reinterpret changed fields.

An operator with a publicly verifiable DNS name publishes the catalog at
\url{https://publisher.example/.well-known/ai-catalog.json}, replacing the
example host with its FQDN. Each resource uses the ARD
domain-anchored identifier form \texttt{urn:air:<publisher-fqdn>:\ldots} and an exact
value-or-reference and media-type representation. An entry can hand a client to MCP, A2A,
OpenAPI, a nested catalog or a versioned TOS Agentic Operation descriptor. The ARD catalog is
stable discovery metadata; rapidly changing price, free capacity, queue depth, local safety
state, model residency and availability are returned only by a fresh operation quote and admission
flow.

TOS can run a separately deployable ARD Registry. It implements the mandatory versioned HTTP
REST search baseline, including \texttt{POST /search}, and may ingest public catalogs,
operator-approved private catalogs, upstream registries, TOS on-chain discovery seeds and
signed profile manifests. Registry responses preserve whether each field came from the
publisher, on-chain state, live observation, attestation or registry-derived ranking.
Federation is plural: no single TOS Registry is mandatory, and referral depth, cycles,
fan-out, time, response size, caching and retries are bounded.

ARD stops before invocation. It does not define TOS authentication, authorization, scheduling,
execution, payment, receipts, evidence or settlement. Before transferring value or invoking
a service, a client independently verifies publisher binding, the current TOS operation descriptor,
runtime and endpoint authorization, manifest commitments, chain activity, quote expiry,
admission and payment destination.

\subsection{Capability Discovery}

Capability discovery separates compact on-chain commitments from extensible off-chain
descriptions. A client can inspect an entry's owner, activity, bond and commitment hashes,
then retrieve a signed descriptor and verify its relationship to chain state.

The current per-actor registry does not solve chain-wide search. TOS ARD Registries enumerate
and rank services while retaining the Capability Registry as one source of on-chain seeds and
commitments. Indexed data is a derived view: clients must verify critical publisher, owner,
authorization, capability, activity, endpoint, quote and payment fields against signed,
observed or on-chain state before committing funds.

\subsection{Names and Reachability}

A service may be reached through a human-readable \texttt{name.tos} identity or a raw ADNL
identity. A descriptor may publish RLDP/TOS Sites and optional HTTPS, QUIC, WebSocket or
streaming bindings without changing the identity and authorization model. Ordinary home and
site networks may require an owner-selected relay or reverse tunnel. A relay may forward
encrypted traffic but must not receive owner authority, unrestricted runtime credentials or
settlement control.

A public root \texttt{.tos} registration product, signed multi-endpoint descriptors,
short-lived health data and production relay service are planned. DNS and TOS Sites
primitives already exist, but they are not equivalent to a complete service discovery
marketplace.

A \texttt{name.tos} identity is not automatically a conventional public-DNS trust anchor for
ARD. A public service without its own DNS domain uses an approved HTTPS gateway that binds an
ARD publisher namespace to the separately signed TOS identity, or it participates in a
private Registry with an explicit \texttt{.tos} resolution policy. ARD publisher control,
\texttt{name.tos}, ADNL, account and manifest commitments remain distinct verifiable
bindings.

\subsection{Locality-First Routing Across Mesh Tiers}
\label{sec:mesh-routing}

The mesh tiers introduced in Section~\ref{sec:agent-mesh} are realized through the
discovery and invocation primitives described earlier in this section, not through a
separate escalation protocol. A client or agent runtime performing an ARD search can rank
candidate resources by measured or declared latency, region and privacy policy in addition
to price and capability, and a federated ARD Registry can apply the same ranking when
resolving a query across upstream catalogs. Preferring the nearest capable node before
considering a broader tier is a client- or registry-side policy expressed through existing
search, ranking and federation behavior, not a new consensus primitive.

A single task may fan out across tiers before it settles. An agent running locally may
resolve routine sub-tasks itself, issue independent quotes to one or more nearby specialized
edge nodes for narrow capabilities such as retrieval, transcription or a domain-specific
tool, and reserve a regional- or frontier-tier service call for the portion of the task that
requires it:

\begin{center}
\texttt{User -> Nearest Edge Agent -> Nearby Specialized Agents -> Regional Compute ->
Cloud Frontier Model}
\end{center}

Each hop in this pattern is an ordinary TOS-authorized invocation: a resolved descriptor, a
live quote, a bounded payment authorization and a signed receipt. ARD discovery lets an
agent find a candidate at any tier; it does not grant that candidate execution, spending or
settlement authority, and a locality-first search ranking is evidence for the requester's
own admission and payment decision, not a substitute for it.


\section{Reference Agentic Operation Profiles}
\label{sec:profiles}

\subsection{Messenger and Mailbox}

Communication is a useful first profile because transport success is not enough. An agent
must know who may contact whom, under what recipient policy, through which relay, with what
retention, delivery attempts, acknowledgement and economic protection against spam.

Representative operations include:

\begin{itemize}
  \item \opid{messenger.send@1};
  \item \opid{messenger.introduce@1};
  \item \opid{mailbox.deposit@1};
  \item \opid{mailbox.read@1};
  \item \opid{mailbox.delete@1}; and
  \item \opid{messenger.acknowledge@1}.
\end{itemize}

The encrypted message body remains off-chain. A request binds sender identity, target,
capability, ciphertext digest, byte and retention limits, relay quote, nonce, expiry and
receipt policy. A sender may deliver directly when the target is online or deposit ciphertext
into one or more authorized Mailboxes.

A typical offline path is:

\begin{center}
\texttt{Sender -> Relay Quote -> Mailbox Deposit -> Stored Receipt}\\
$\longrightarrow$ \texttt{Target Read -> Delivery Receipt -> Acknowledgement}\\
$\longrightarrow$ \texttt{Settlement or Bond Release}
\end{center}

The receipt levels remain distinct:

\begin{itemize}
  \item \textbf{Stored receipt:} an authorized Mailbox accepted ciphertext under a declared
        retention and replica policy.
  \item \textbf{Delivery receipt:} a target device or relay accepted the encrypted item.
  \item \textbf{Acknowledgement receipt:} the target agent emitted an application-level
        acknowledgement.
  \item \textbf{Settlement receipt:} payment, refund or bond state reached the declared
        finality.
\end{itemize}

A Mailbox cost can be modeled as

\[
  C_{\mathrm{mailbox}} =
  C_{\mathrm{admission}}
  + p_b B T
  + p_r R
  + p_a A
  + C_{\mathrm{priority}},
\]

where $B$ is ciphertext bytes, $T$ retention duration, $R$ replica count and $A$ delivery
attempts. Implementations may use prepaid allowance rather than settling every deposit
individually.

For known contacts, a recipient may grant a zero-price capability with a daily quota. For an
unknown sender, \opid{messenger.introduce@1} may require a refundable attention bond. If the
recipient accepts or replies, the bond can be returned; if policy determines the request was
abusive, part may be paid to the recipient or forfeited under a disclosed rule. Pricing alone
is not sufficient: capabilities, quotas, reputation, congestion controls and bounded fan-out
remain mandatory.

Current Messenger and Mailbox code provides profile-specific foundations such as signed
capabilities, nonce and expiry checks, deposit/read/delete operations and stored
acknowledgements. A frozen generic communication operation specification, commercial relay
lease, inbox bond and cross-provider conformance suite remain partial or planned.

\subsection{Model and Tool Invocation}

Model and tool calls demonstrate multidimensional metering and evidence. Representative
operations include:

\begin{itemize}
  \item \opid{model.infer@1};
  \item \opid{embedding.create@1};
  \item \opid{speech.transcribe@1};
  \item \opid{tool.invoke@1}; and
  \item \opid{workflow.execute@1}.
\end{itemize}

A model operation binds an exact model or service revision, input and output schemas, context
and output limits, latency class, region, privacy and egress policy, tool permissions,
evidence level and maximum cost. Resource dimensions may include input tokens, output tokens,
accelerator milliseconds, cached context, tool calls and external provider fees.

A provider advertises a bounded service, not raw GPU access or arbitrary consumer-supplied
code. The provider admits the request under local policy, executes through an approved adapter
and returns a usage and result receipt. The receipt may state a provider signature,
standardized benchmark, replicated execution, hardware attestation or stronger proof, but the
evidence level must not be overstated.

A model result receipt does not reveal private prompts or full output by default. It binds
content digests and any evidence references permitted by the data policy.

\subsection{Storage and Event Delivery}

Storage operations include:

\begin{itemize}
  \item \opid{storage.put@1};
  \item \opid{storage.get@1};
  \item \opid{storage.renew@1};
  \item \opid{storage.delete@1};
  \item \opid{event.publish@1}; and
  \item \opid{event.subscribe@1}.
\end{itemize}

Storage profiles meter bytes, retention duration, durability tier, replica count, retrieval
bandwidth and region. A lease receipt identifies the content digest, admitted duration and
durability claim. Deletion semantics distinguish removal from one provider, expiration of a
lease and cryptographic erasure where supported.

Event profiles meter event count, payload bytes, subscriber fan-out, retention and replay.
Subscriptions carry a bounded filter, rate and backlog policy. A publisher cannot create
unbounded subscriber work, and a disconnected subscriber cannot force unlimited durable
queue growth.

\subsection{Commerce, Booking and Human Services}

Commerce and reservation profiles demonstrate that Agentic Operations are not limited to
digital compute.

Representative operations may include:

\begin{itemize}
  \item \opid{commerce.quote@1};
  \item \opid{commerce.order@1};
  \item \opid{booking.reserve@1};
  \item \opid{booking.cancel@1};
  \item \opid{delivery.request@1}; and
  \item \opid{human.task@1}.
\end{itemize}

These profiles require domain-specific state machines. A booking operation defines resource
identity, dates or duration, capacity, price, cancellation, deposit and confirmation. A
commerce order defines item revision, quantity, fulfillment, tax or jurisdictional policy and
refund conditions. A human task defines scope, deadline, compensation, privacy, evidence and
dispute policy.

TOS supplies common identity, authority, budget, receipt and settlement primitives; it does
not replace consumer law, tax, sanctions, labor, licensing or data-protection obligations that
may apply to the operator.

\subsection{Task and Escrow Composition}

Task Escrow can compose multiple Agentic Operations. A creator funds a task, an assigned
agent or open claimant accepts it, and the agent invokes communication, search, compute,
storage or human services under delegated capability and sub-budgets. The final task result
references the receipts of those operations.

A task budget is not permission for every subordinate action. The controller, operation,
target, cumulative spend, deadline and evidence policy remain bounded. Unused sub-budgets are
returned according to the task and settlement rules.
\section{Physical Edge Intelligence Profile}

\subsection{Terminal as the Product Unit}

The TOS AI Edge Computing Terminal is one planned owner-operated execution profile. It
combines measured local resources, approved runtime adapters, bounded task admission, TOS
identity and reachability, metering and signed receipts. A terminal may be a Linux GPU
workstation, an AI PC, a small edge server or an industrial ARM/Jetson appliance.

Consumers schedule against a service capability such as text generation, embedding, OCR,
speech recognition or video analysis. A capability states the exact model or task revision,
schemas, limits, privacy and region policy, evidence level, price and expiry. Hardware model
numbers and vendor performance claims are evidence used by the scheduler; the consumer does
not receive raw access to the device.

Consistent with the deployment discussion in Section~\ref{sec:model-scale-shift}, the
terminal protocol is model-size and vendor neutral. A reference implementation may initially
target model classes that fit owner-operated workstations, AI PCs and Jetson-class appliances,
but no operation identifier, receipt or settlement rule depends on a parameter count,
accelerator or model family. Smaller and larger models are valid when an approved runtime
adapter and bounded resource envelope support them.

The set of eligible terminal hardware is intentionally broad within that class: a Mac
mini, MacBook, gaming PC, home server, NAS appliance or small industrial workstation can
register under the same Capability Registry and Service Actor primitives described earlier
in this paper as a rack-mounted GPU server, provided it passes the same admission,
isolation and evidence requirements described in this section. What a terminal contributes
is not limited to raw inference: a Service Actor entry can equally represent local storage,
request routing to another known-good node, retrieval or context-memory services built on
stored state, or a narrowly scoped tool integration. The product requirement is a bounded,
evidenced, receipted capability, not a particular hardware class or deployment scale.

\subsection{Runtime and Executor Isolation}

The public service listener is separated from the administrative interface. Public adapters
run under dedicated operating-system identities, containers or equivalent sandboxes with
explicit CPU, RAM, accelerator-memory, filesystem, network and tool grants. They receive no
Docker socket, shell, owner key, unrestricted wallet key, home directory or raw physical-I/O
authority.

The operator approves every model and executable artifact in the initial product. A runtime
adapter performs compatibility checks, bounded load and unload, cancellation, error
normalization, usage reporting and cleanup after success, timeout, disconnect, out-of-memory
failure or process crash. Consumer-supplied containers, native programs and models are not
accepted by the TOS terminal product.

\subsection{Admission, Scheduling and Resource Bounds}

The local scheduler is authoritative for admission. It considers model compatibility, RAM,
accelerator memory, storage, queue and in-flight limits, context and output bounds, deadlines,
owner reservations, thermal policy, price, privacy and evidence requirements. Remote ARD
discovery cannot reserve local hardware merely by publishing an optimistic capacity value.
A physical-terminal catalog advertises a bounded site or fleet-broker capability, not raw
per-device administration, site topology, sensor streams or actuator handles.

Every implementation must bound connections, sessions, streams, nonces, quotes, queues,
retries, model caches, temporary files, logs, receipt queues, chain watchers and durable
journals. Completion, cancellation, failed payment, adapter crash, upgrade and shutdown each
need a tested cleanup path. These constraints are required to prevent anonymous RAM, VRAM,
disk, watcher and durable-state growth.

\subsection{Site-Bound Physical AI}

A physical terminal runs beside sensors, robots, vehicles or industrial equipment. Its
mandatory priority order is:

\begin{enumerate}
  \item emergency and independent safety interlocks;
  \item deterministic control deadlines;
  \item local real-time perception and sensor fusion;
  \item local asynchronous analysis and maintenance;
  \item approved external service requests;
  \item background update, telemetry and compaction.
\end{enumerate}

TOS does not enter the hard real-time loop. Raw CAN, GPIO, serial, fieldbus, camera or actuator
interfaces are never public capabilities. A narrow semantic action must pass local policy and
an independent safety controller, which may reject an otherwise valid network request.

\subsection{Disconnected Operation and Reconciliation}

A physical terminal must continue its approved local workload while disconnected. It may use
only unexpired cached policy and explicitly bounded offline authority. Events and obligations
are written to a size- and age-bounded, tamper-evident journal. Reconnection first observes
revocation and policy expiry, then uploads and reconciles events idempotently. A reconnect
must not execute a paid task twice, lose a refund obligation or grow the journal without
limit.

\subsection{Safe Model and Software Updates}

Model, runtime, firmware and policy updates are content-addressed, signed by an authorized
release identity, compatibility-checked and staged before activation. Terminals maintain an
active slot and a known-good rollback slot. Activation is crash-safe and subject to health
gates, anti-rollback policy and bounded artifact retention. Fleet rollouts use canary and
progressive rings and stop automatically when declared health thresholds fail.

\subsection{Fleet Management}

Fleet state is primarily off-chain. A fleet service handles enrollment, scoped delegation,
site grouping, staged rollout, health, revocation, offline expiry and retirement without
publishing detailed topology or continuous telemetry on-chain. Commands have bounded fan-out,
expiry and idempotency. Compromise of a fleet controller must not confer owner custody or
override independent local safety policy.


\section{Deployment Profiles and Market Thesis}

\subsection{Model Scale and the Deployment Shift}
\label{sec:model-scale-shift}

TOS is model-size neutral. No protocol rule depends on a particular parameter count,
accelerator vendor or foundation-model family. Current deployment practice nevertheless spans
several broad tiers:

\begin{longtable}{L{0.18\textwidth}L{0.32\textwidth}L{0.39\textwidth}}
\toprule
Deployment tier & Representative location & Typical operation profile \\
\midrule
\endhead
Constrained local &
Phones, wearables, embedded and IoT devices &
Voice, local policy, sensor interpretation and lightweight tools \\
Owner-operated edge &
Workstations, AI PCs, small servers and industrial appliances &
Agent runtimes, retrieval, tool use, private inference, storage and local automation \\
Regional or enterprise &
Enterprise servers and aggregation infrastructure &
Larger models, knowledge systems, replicated services and organization policy \\
Cloud frontier &
Large data centers &
Frontier training and the heaviest reasoning or multimodal workloads \\
\bottomrule
\end{longtable}

Model and hardware bands change faster than the operation protocol. Provider descriptors and
quotes state actual capability and limits; the whitepaper does not freeze a preferred model
size. The durable thesis is that agentic work will be distributed across independently
operated local, edge, regional and cloud resources.

\subsection{From Cloud-Centric AI to an Agent Mesh}
\label{sec:agent-mesh}

A cloud-centric pattern sends every request directly to one remote provider:

\begin{center}
\texttt{User or Agent -> Centralized Service}
\end{center}

An agent mesh can resolve a request through progressively broader resources:

\begin{center}
\texttt{Local Agent -> Nearby Provider -> Regional Service -> Cloud Frontier Service}
\end{center}

A single task may use several tiers. A local agent can handle policy and context, invoke a
nearby private service for one operation, and escalate only the narrow subtask that requires a
regional or frontier provider. Each external step remains an ordinary Agentic Service
Operation with an exact revision, quote, bounded authorization and receipt.

TOS does not decide which tier is correct. Discovery clients and provider policy rank
candidates by capability, measured latency, price, privacy, evidence and locality. TOS makes
the resulting action accountable across those tiers.

\section{Receipts, Evidence and Verifiable Workflows}

A receipt is the bridge between external action and shared economic state. TOS does not claim
that every service result can be verified by consensus. It standardizes what was requested,
who accepted it, what resources were reported, what result commitment was returned and which
evidence and dispute policy the parties accepted.

\subsection{State and Evidence Boundary}

TOS records what consensus can enforce:

\begin{itemize}
  \item account authority and spending limits;
  \item operation, quote, capability and budget commitments;
  \item task assignment, status, budget and deadlines;
  \item receipt, settlement and dispute authority;
  \item result, evidence and metadata commitments;
  \item registry bond, activity and reputation references.
\end{itemize}

The chain should not store encrypted message bodies, private prompts, service credentials,
full model transcripts, large files or raw sensor streams. These remain in content-addressed
storage, encrypted communication systems or provider infrastructure.

\subsection{Proof Adapters}

Verification may use signatures, attestations, proof-backed transport evidence, committee
decisions or domain-specific proof systems. A proof adapter translates external evidence into
a compact decision or commitment understood by a task or verifier actor. The core protocol
remains neutral to the model provider and proof backend.

Claims must state their evidence level. A terminal declaration, an observed successful call,
a standardized benchmark, an independent audit, hardware attestation and replicated execution
are different forms of evidence and must not be presented as interchangeable. A signed
receipt proves that a signer made a statement about an operation; its semantic meaning still
depends on the accepted policy and trust model.

\subsection{Generic Agentic Operation Workflow}

\begin{enumerate}
  \item A client resolves a provider identity and verifies its signed operation descriptor.
  \item The client requests a quote for an exact operation revision, resource limits and
        evidence policy.
  \item A controller, user or sponsor authorizes bounded resource use and maximum payment.
  \item The provider admits the request only after capability, local policy and resource
        checks.
  \item An isolated adapter, relay or service executes the accepted operation.
  \item The provider signs a result and usage receipt linked to the operation identity.
  \item The client verifies the receipt and the settlement path releases, nets or refunds
        value.
\end{enumerate}

Every step is asynchronous. Failure, timeout and dispute are explicit workflow states rather
than hidden exceptions in an off-chain orchestration process.

\subsection{Physical-Terminal Workflow}

A site terminal receives sensor data locally, schedules safety and real-time perception ahead
of all network work, applies local policy and emits a bounded event or service result. If the
site is disconnected, it continues approved local operation and journals the event. After
reconnection it observes revocations, reconciles the journal idempotently and publishes only
the receipt or evidence allowed by site data policy. Blockchain availability is therefore
not a prerequisite for an emergency stop or deterministic local control decision.


\subsection{Messenger and Mailbox Workflow}

\begin{enumerate}
  \item A sender resolves the recipient policy and one or more authorized relay or Mailbox
        descriptors.
  \item The sender obtains a quote or uses an existing allowance and signs a bounded
        \opid{messenger.send@1} or \opid{mailbox.deposit@1} request.
  \item The relay verifies signature, nonce, expiry, capability, byte limit, quota and budget
        before allocating storage or delivery work.
  \item The encrypted payload is stored or delivered off-chain.
  \item The Mailbox returns a stored receipt; the recipient device may later return a delivery
        or acknowledgement receipt.
  \item Payment, refund or attention bond is settled according to the accepted policy.
\end{enumerate}

A stored receipt does not reveal message plaintext and does not prove the recipient read or
approved the content.

\subsection{Model and Tool Workflow}

\begin{enumerate}
  \item The caller selects an exact model or tool operation revision and provider.
  \item The quote fixes token, time, tool and external-cost dimensions and a maximum charge.
  \item The provider verifies capability and admits the request under local resource policy.
  \item The approved runtime executes with bounded context, output, time, tools and egress.
  \item The provider returns a result digest, actual usage, charge and declared evidence.
  \item The caller verifies the receipt and settles, refunds or disputes.
\end{enumerate}
\section{APIs, SDKs and Operator Tooling}

\subsection{tosctl}

\texttt{tosctl} is the canonical operator tool for node configuration, native wallets and agent
actor workflows. The implemented command families create and inspect Agent Wallet profiles,
deploy and manage Agent Accounts, create and operate Task Escrows, and deploy or update
Capability Registry entries.

Human-readable output supports operators; JSON output supports agent runners and automation.
Secrets remain in the configured vault rather than the profile JSON.

\subsection{Read APIs}

The authenticated node-control service exposes read endpoints without URI version prefixes:

\begin{itemize}
  \item \texttt{GET /agents/\{address\}};
  \item \texttt{GET /tasks/\{address\}};
  \item \texttt{GET /tasks} with status, creator, agent, deadline and pagination filters.
\end{itemize}

Task-list chain reads use bounded concurrency, deterministic ordering, individual timeouts and
failure isolation. The list enumerates tasks registered in the local node-control
configuration; it is not a complete chain-wide index.

Public API errors distinguish invalid requests, missing state, unavailable RPC, invalid
contract state and timeout without exposing raw internal transport errors.

\subsection{Trust Tiers}

An agent controlling material funds should verify state through its own node or a proof-backed
client. A trusted remote API may be suitable for an operator's own infrastructure. Third-party
indexers are useful for search and analytics but must not become authority for balances,
permissions or settlement.


\subsection{Agent and Provider SDKs}

The common operation protocol requires at least:

\begin{itemize}
  \item an Agent SDK for discovery, quote comparison, capability use, budgeting, invocation,
        receipt verification and settlement;
  \item a Provider SDK for descriptor publication, admission, metering, receipt signing,
        refund and restart recovery;
  \item a Receipt Verifier library independent of the provider runtime;
  \item operation-definition tooling for schema and metering manifests; and
  \item profile conformance runners with canonical success and failure vectors.
\end{itemize}

At least two independently authored implementations should interoperate for a profile before
that profile is described as open and portable.

\section{Security and Anti-Abuse}
\label{sec:anti-abuse}

\subsection{Economic Denial of Service and Spam Resistance}

Pricing is one anti-abuse mechanism, not the only one. A wealthy attacker may still pay to
create harmful load, and a legitimate trusted relationship may reasonably use a zero-price
allowance. A production operation profile combines:

\begin{itemize}
  \item authenticated identity and replay-resistant requests;
  \item narrowly scoped capabilities;
  \item per-caller, per-target and per-capability quotas;
  \item maximum resource and cost envelopes;
  \item refundable bonds where recipient attention or scarce reservation is imposed;
  \item congestion-dependent pricing and admission;
  \item reputation-sensitive policy where the profile can support it safely;
  \item bounded fan-out, retries, backlog and retention;
  \item queue, memory, disk and durable-state limits; and
  \item provider- and recipient-defined local policy.
\end{itemize}

The initiator should fund or hold an allowance for the external cost it imposes. The purpose
is not to make human or agent communication expensive. The purpose is to prevent an identity
from making storage, compute, delivery, attention or retry costs effectively infinite for
someone else.

\subsection{Cheap Checks Before Expensive Work}

A public provider processes untrusted requests in the following order:

\begin{enumerate}
  \item enforce framing, byte and schema limits;
  \item verify signature and domain binding;
  \item verify nonce, expiry and idempotency identity;
  \item verify capability, revocation and delegation;
  \item verify quota, quote and maximum budget;
  \item reserve bounded local resources or payment; and
  \item only then decrypt large payloads, allocate storage, load a model, call an external
        tool or notify a recipient.
\end{enumerate}

Failure at an early step releases temporary state. An invalid request must not trigger
expensive parsing, model loading, attachment expansion, network fan-out or durable journal
growth.

\subsection{Attention Bonds and Recipient Policy}

Some operations consume a scarce human or agent resource: attention. A communication profile
may distinguish established contacts from introductions.

An established contact can receive a zero-price, quota-limited capability. An unknown sender
may lock a refundable bond. Acceptance or reply can release the bond; rejection under a
declared anti-spam policy can forfeit a bounded part. The recipient cannot arbitrarily change
the rule after delivery, and the sender knows the maximum loss before authorization.

Attention bonds are optional profile policy, not a universal tax on communication. They are
one tool alongside allowlists, proof of relationship, organization sponsorship and local
filtering.

\subsection{Amplification and Autonomous Loops}

Agent systems can amplify work through delegation, subscriptions, retries and event fan-out.
Every operation carries or derives:

\begin{itemize}
  \item maximum delegation depth;
  \item maximum child-operation count;
  \item cumulative budget and resource ceilings;
  \item maximum retry count and backoff;
  \item maximum recipients, subscribers or replicas; and
  \item a cancellation path that propagates through the causal chain.
\end{itemize}

A parent operation may allocate sub-budgets, but children cannot exceed the parent's
authorized total. Cycles in workflows, federation and referral graphs are detected or bounded.
\subsection{Key Compromise}

Owner keys belong in operator-controlled vaults. Controller keys belong to runtimes only when
their on-chain authority is acceptably bounded. Rotation must be possible without changing
the owner. Runtime manifests must never contain owner secrets.

\subsection{Replay and Domain Confusion}

Signed actions require sequence numbers and expiry. Protocols should bind signatures and
permissions to the intended network, account, task and operation. A valid message for one
workflow must not authorize another.

\subsection{Escrow Safety}

Contracts reject out-of-order messages, unauthorized senders and payouts above available
balance. Deadlines, review windows and terminal states are explicit. Local CLI validation is
not a substitute for contract checks.

\subsection{Service and Evidence Risk}

A capability claim is not proof of service quality. A metadata hash is not proof that its
content is true. A result hash is not proof that a result is correct. Bonds, verifier roles,
attestations and reputation can reduce risk, but task settlement policy must state which
evidence is authoritative.

\subsection{Executor and Supply-Chain Isolation}

Public requests are untrusted input. Runtime adapters must enforce schema, byte, dimension,
duration and output limits before allocating expensive resources. Models, runtimes and
updates require authenticated provenance, canonical hashes, compatibility checks and
operator approval. Sandboxing, least-privilege device access, network egress policy and
secret separation limit the impact of a compromised model server or parser. The public
service interface must never become an administrative or arbitrary-code execution path.

\subsection{Physical Safety}

Possession of a valid TOS key, payment authorization or network receipt does not grant raw
actuator authority. Local semantic capability policy and independent safety interlocks remain
authoritative. Safety and control deadlines pre-empt external work; overload or network
failure must fail toward a locally defined safe state.

\subsection{Offline Authority and Update Risk}

Offline operation creates stale-policy and replay risk. Offline rights therefore need explicit
scope, expiry and budgets, while journals need integrity, bounded retention and idempotent
reconciliation. Update systems require separate release authority, anti-rollback checks,
known-good recovery and staged fleet rollout. A terminal that cannot verify a package or
recover a healthy runtime must continue only its previously approved safe workload.

\subsection{Fleet Compromise}

Fleet controllers use scoped delegation and revocable terminal identities. Commands carry
target scope, expiry and idempotency identifiers. Rollout fan-out, retry queues, health
history and offline inventory are bounded. Fleet compromise must not reveal owner custody
keys, customer payloads or unrestricted site topology.

\subsection{ARD Catalog and Registry Risk}

Catalogs, descriptions, representative queries, endpoints, nested references and federated
responses are untrusted input. Implementations validate the pinned ARD schema, publisher
domain, resource identifiers, trust metadata, value-or-reference rules and media types.
Natural-language fields are data, never instructions to an agent. Crawlers allow only
approved schemes, ports, redirects and address classes; they defend against DNS rebinding,
server-side request forgery, oversized or compressed input, recursive catalogs, federation
cycles, stale rollback, equivocation and ranking manipulation.

Catalog, field, entry, reference, redirect, hop, response, time, retry, cache, index and
per-publisher limits prevent anonymous memory, disk and work growth. Visibility policy
prevents private endpoints, credentials, topology, sensor data and customer payloads from
entering public search or federation. Registry output retains field provenance and cannot
grant wallet, terminal, model, update, fleet or actuator authority.

\subsection{Indexer Authority Drift}

Search results and reconstructed timelines are non-authoritative. Critical clients should
re-read relevant contract state before transferring funds, accepting work or resolving a
dispute.

\subsection{Availability and Message Amplification}

Automated actors can create unbounded retry loops and fan-out. Production workflows require
funded or quota-backed bounded retries, backoff, timeout policy, cumulative child-operation
budgets, queue limits and observable failure records.
Terminals additionally require soak tests for RAM, accelerator memory, disk, watchers,
connections, offline journals and update artifacts across success and every failure path.

\section{Protocol Economics and Token Launch Discipline}

\subsection{Product Before Emission}

The native asset can coordinate validator security, settlement and open participation, but it
cannot create product-market fit. Token incentives launched before repeat organic demand can
distort customer discovery, attract mercenary volume and make economic parameters harder to
change. TOS therefore treats operation adoption, sustainable unit economics and independent
provider participation as prerequisites for activating discretionary service-linked issuance.

The practical sequence is:

\begin{enumerate}
  \item ship the compute reference profile and use it in a first-party application;
  \item win unrelated design partners and repeat payers without AIPoW rewards;
  \item demonstrate cross-provider conformance, low-cost settlement and non-subsidized
        retention;
  \item publish the decentralization, operating-funding and governance path; and
  \item only then evaluate AIPoW shadow data and any proposal for activation.
\end{enumerate}

This sequencing follows a general crypto product discipline: a token may amplify an existing
network, but it should not substitute for a product users retain.
\footnote{Miles Jennings and Jason Rosenthal, ``Getting ready to launch a token: What you need
to know,'' a16z crypto, April 25, 2024:
\url{https://a16zcrypto.com/posts/article/getting-ready-to-launch-a-token}.}
No token allocation, reward schedule or valuation narrative is evidence of demand.

\subsection{Native Unit, Utility and Economic Character}

TOS is the native settlement and security asset of TOS Network. One TOS equals $10^9$
\nanotos, the smallest indivisible unit. Contract balances, fees, task budgets, bonds,
spending limits and validator stakes are represented in \nanotos{} to avoid rounding
ambiguity.

TOS is designed for functional use within the network:

\begin{itemize}
  \item paying transaction, message, storage and Agentic Service Operation fees;
  \item bonding elected validators and securing consensus;
  \item compensating elected validators for producing and validating blocks;
  \item funding task escrow, refunds, operation settlement and refundable anti-abuse bonds;
  \item bonding capability, verifier and other protocol roles where specified; and
  \item participating in protocol governance only where an accepted protocol rule grants
        that function.
\end{itemize}

TOS does not represent equity, debt, a partnership interest or ownership of a legal entity.
It gives no contractual claim on project assets, revenue, profits, dividends or cash flows;
no right of redemption by a company, foundation or operator; no guaranteed return, fixed
yield, liquidity or price support; and no ownership of a model, terminal, data set or
off-chain service provider. Holding TOS alone does not produce a protocol reward. Validator
compensation requires election, locked stake, active operation and continuing exposure to
the protocol's complaint and penalty rules.

\subsection{Supply Policy}

The initial distribution is designed to avoid a large insider premine. Approximately
500 million TOS is created through ordinary validator block rewards; the remaining
community-agent allocation (approximately 4.5 billion TOS of the five-billion total-supply
policy) is created through Artificial Intelligence Proof of Work (AIPoW), a separate protocol reward
mechanism described below, is not funded at genesis and is never held by a treasury
wallet. The current production policy is:

\begin{longtable}{L{0.48\textwidth}L{0.38\textwidth}}
\toprule
Item & Policy target \\
\midrule
\endhead
Smallest unit & 1 \nanotos \\
Unit conversion & 1 TOS = 1,000,000,000 \nanotos \\
Approximate total network supply target & 5,000,000,000 TOS \\
Approximate validator-reward creation target & 500,000,000 TOS \\
Community-agent allocation (Artificial Intelligence Proof of Work) & 4,500,000,000 TOS \\
Provisional main-wallet genesis balance & 100,000 TOS \\
Elector operating reserve & 500 TOS \\
Configuration-contract operating reserve & 500 TOS \\
Provisional total genesis creation & 101,000 TOS \\
Reference post-genesis validator creation & 499,899,000 TOS \\
Proof-of-work giver allocation & 0 TOS \\
Team, investor, foundation and treasury allocation & 0 TOS \\
Target validator distribution duration & Approximately seven years \\
Treatment of network outages & No creation and no later catch-up \\
Reward termination & Configuration governance tapers or sets ConfigParam 14 to zero \\
New hard consensus supply cap & None in this policy \\
\bottomrule
\end{longtable}

Five hundred million TOS of validator creation, five billion TOS of total supply and seven
years are policy calibration targets, not an exact consensus cap, a guaranteed completion
date or a statement about future market value. The actual total and completion date depend
on finalized block production, shard behavior, validator availability, configuration
activation and the eventual governance transaction that stops native block creation. The
community-agent allocation is created only through Artificial Intelligence Proof of Work, described in
its own subsection below, with its own launch gates; it is not distributed through
ConfigParam 14 or the Elector.

\subsection{Genesis Bootstrap and Allocation}

The production zerostate provisionally creates 101,000 TOS: a bounded 100,000-TOS
validator-bootstrap main wallet and 500-TOS operating reserves for each of the Elector and
Configuration contracts. It contains no proof-of-work giver, faucet, discretionary mint
contract, general treasury, market-making inventory or native allocation to a team,
investor, foundation or ecosystem fund. The protocol distribution plan includes no sale of
a large genesis inventory.

The zerostate commits four original validators with unique signing and network identities
and equal initial consensus weight. The main wallet may transfer to each validator's
published controlling masterchain wallet exactly 20,000 TOS of principal for two overlapping
minimum-stake elections, plus no more than 100 TOS of separately measured wallet,
message and retry costs. The signing keys, controlling wallets, allowed destinations,
transfer caps and transactions must be public.

After two overlapping elected validator sets have been installed successfully, all remaining
main-wallet TOS must be transferred to a published unspendable address, leaving the wallet
with zero spendable balance. The main wallet is a temporary bootstrap mechanism and must not
become a treasury, fee collector, market-making account, governance fund or emergency mint
authority. The two system-contract reserves are controlled only by their deployed contract
code and are not validator rewards or discretionary allocations.

\subsection{Validator-Led Native Creation}

Post-genesis native creation uses the existing block-value and Elector path:

\begin{center}
\texttt{ConfigParam 14 block creation}\\
$\longrightarrow$ \texttt{configured collector or Elector fallback}\\
$\longrightarrow$ \texttt{active validator-set bonus pool}\\
$\longrightarrow$ \texttt{stake and bonus recovery}
\end{center}

For validator $i$, the existing proportional rule is

\[
  B_i =
  \left\lfloor
    B_{\mathrm{set}}\frac{E_i}{\sum_j E_j}
  \right\rfloor,
\]

where $B_{\mathrm{set}}$ is the active set's recoverable bonus pool and $E_i$ is
validator $i$'s effective stake. Deterministic integer-remainder handling, stake returns,
complaints and penalties follow the existing Elector rules. There is no separate proposer
allocation, passive-holder yield, work-score reward, vote-count reward, private mint key or
off-chain reward discretion.

Penalties are proportional to the stake a validator controls rather than a flat amount. The
misbehaviour schedule in \texttt{ConfigParam 40} sets a base fine that the severity and
duration of the fault scale up, with the most serious tier reaching TM\$2{,}500 plus one
quarter of the stake. Duration thresholds are expressed as fractions of the validation round
so that every tier is reachable within one. Two consequences are worth stating: an operator
that gathers more stake carries proportionally more at risk rather than the same fixed
exposure, and any contract holding stake on behalf of others can size the operator's own
required funds from this schedule instead of guessing. A network that omits the parameter
falls back to a flat fine with no proportional component, which makes the amount at risk
independent of the amount controlled.

A validator must submit a valid bid, be selected, keep its stake locked, operate its assigned
consensus responsibilities and remain eligible under the protocol rules. Native block
creation and collected transaction fees are different economic flows: only ConfigParam 14
creation increases gross native supply, while fees transfer existing TOS.

\subsection{Approximate Seven-Year Distribution}

The reference post-genesis validator creation target is 499,899,000 TOS. For calibration,
the expected creation rate is modeled as

\[
  \dot S =
  \lambda_{\mathrm{mc}}R_{\mathrm{mc}} +
  \sum_s \lambda_s\frac{R_{\mathrm{bc}}}{2^{d_s}},
\]

where $\lambda_{\mathrm{mc}}$ is the measured finalized masterchain block rate,
$R_{\mathrm{mc}}$ is the configured masterchain creation value, $\lambda_s$ is the
finalized rate of basechain shard $s$, $d_s$ is its prefix depth and $R_{\mathrm{bc}}$ is
the configured basechain value. Depth adjustment is intended to keep aggregate basechain
creation approximately stable across balanced shard splits, but actual finalized chain data
remains authoritative.

The current implementation candidate uses 0.569879384 TOS per masterchain block and
0.335223167 TOS per unsplit-basechain block. At the locally measured reference rate of
approximately 2.5 finalized blocks per second for each chain, those values project nearly
the full post-genesis target from August 1, 2026 to August 1, 2033. They are provisional
until sustained multi-validator, multi-host, outage and shard split-and-merge measurements
pass the published launch gates.

Issuance is event-based:

\begin{center}
\texttt{no finalized produced block} $\Longrightarrow$ \texttt{no native creation}.
\end{center}

An outage, partition, failed election or delayed restart creates no emission debt. There is
no backfill, catch-up multiplier, anniversary tranche, pending-emission queue or automatic
extension of a seven-year calendar. Faster block production can reach the target earlier;
slower production can reach it later. These outcomes must be measured and disclosed rather
than described as protocol faults.

Governance must monitor finalized creation and begin a public taper review before projected
gross validator creation reaches 495 million TOS. It may reduce ConfigParam 14 to manage
activation delay and must set both masterchain and basechain creation values to zero near
the 500-million validator-creation policy target. Transaction and service fees continue
after native creation stops. Any later proposal to restart native creation is a separate monetary-policy decision
requiring prominent public review.

\subsection{Artificial Intelligence Proof of Work: The Community Distribution}
\label{sec:aipow}

The community-agent allocation of approximately 4.5 billion TOS is created through
\textbf{Artificial Intelligence Proof of Work (AIPoW)}. The useful-work category is not new:
Bittensor pays participants for validator-scored digital commodities,
Gensyn focuses on execution and verification of machine-learning work, and Render operates a
distributed GPU-work marketplace.\footnote{Primary project descriptions:
\url{https://www.bittensor.com/docs},
\url{https://docs.gensyn.ai/the-gensyn-protocol}, and
\url{https://know.rendernetwork.com/}.} TOS does not claim that attaching a token to AI
compute is itself an innovation. Its narrower design claim is that eligible compute, storage, verification and
physical-edge Agentic Service Operations can share one outcome ledger, and that bounded native
issuance can be derived from unrelated, evidence-graded, on-chain settled demand rather than
from declared capacity or a subjective popularity vote.

The 4.5-billion figure is a provisional policy target, not a launch entitlement. Before any
activation, governance must revisit the emission coefficient, epoch ceiling, cold-start floor
and cumulative cap using observed organic demand, provider margins, network-security needs and
concentration data. Monetary parameters should evolve from measured network use; preserving a
headline allocation is not a reason to issue tokens into weak or circular demand.

AIPoW borrows four economic properties from proof-of-work issuance: permissionless earning,
protocol-automatic creation, no reward for merely holding a token, and no treasury that first
receives the allocation. It does \emph{not} replace Bitcoin's consensus work. TOS consensus is
secured by stake-bonded validators; AIPoW provisions the service-capacity layer. This separation
avoids asking non-deterministic useful work to provide block-timing, fork-choice and inexpensive
universal-verification properties for which ordinary hash puzzles were designed. Those are
known hard requirements when useful computation is proposed as consensus work.\footnote{A.
Merlina, T. Garrett and R. Vitenberg, ``On Replacing Cryptopuzzles with Useful Computation in
Blockchain Proof-of-Work Protocols,'' 2024: \url{https://arxiv.org/abs/2404.15735}.} The name
AIPoW is therefore an issuance and service-provision analogy, not a claim that AI work secures
TOS consensus and not a legal classification of any participant relationship.

\subsubsection{Eligible Work and Evidence-Weighted Score}

Let $J_e$ be the set of unique eligible Agentic Service Operations finalized in epoch $e$.
For operation $j$, let $o_j$ be its frozen \texttt{operation\_id} and revision, $p_j$ its
actual settled price in \nanotos, $x_j$ its measured work units under the corresponding
metering-schema hash, $r_{c_j}$ the published maximum rate per unit for capability class
$c_j$, $m_{\ell_j}$ the evidence multiplier in basis points for evidence level $\ell_j$,
and $q_j \in [0,10000]$ a capability-specific quality and service-level factor. Let $I_j$
equal one only when the operation has a unique replay-protected identity, reached final
settlement, remained outside dispute or forfeiture, used the frozen operation revision,
metering schema and rate card, and satisfies the published payer-independence and evidence
rules. Its integer score contribution is

\[
  v_j = I_j\left\lfloor
    \min\!\left(p_j, r_{c_j}x_j\right)
    \frac{m_{\ell_j}}{10^4}
    \frac{q_j}{10^4}
  \right\rfloor .
\]

Self-declared capacity has $m_{\mathrm{declared}}=0$. A paid receipt may establish that an
operation was requested and settled, but it does not by itself prove semantic correctness.
Higher evidence levels require capability-appropriate evidence such as a deterministic
replay, randomized replicated execution, a benchmark with a hidden challenge seed, a hardware
attestation bound to the service revision, an audited measurement, or a cryptographic proof.
An LLM-as-judge score may be one input for an open-ended task but is not sufficient by itself
for the highest evidence tier. Every multiplier and quality function is a bounded rational
integer rule in the published methodology; floating-point or scorer-local behavior cannot
affect an epoch root.

Operation pricing and AIPoW eligibility are independent. A valid paid operation may receive
zero issuance score. Communication, Mailbox retention, broadcast, event fan-out and other
profiles that are easy to self-generate default to zero AIPoW weight unless a separately
audited methodology demonstrates independent demand, resource evidence and resistance to
circular traffic. An operation definition therefore carries no implied reward eligibility.

A production scoring row additionally commits at least the
\texttt{operation\_id}, operation revision, metering-schema hash, actual usage vector, actual
cost, payer control domain, provider control domain, evidence level and receipt-schema
revision.

The price cap prevents a provider from turning an unusually high invoice into unlimited
score, while the settlement cap prevents a rate card from crediting value no customer paid.
It does not, by itself, prove that payer and earner are independent. Organic settled value is
therefore a separate quantity:

\[
  O_e = \sum_{j\in J_e} I^{\mathrm{organic}}_j
        \min\!\left(p_j,r_{c_j}x_j\right),
\]

where $I^{\mathrm{organic}}_j=1$ only for a settled operation whose payer and earner are in
different disclosed and inferred control domains, which is not a protocol-funded challenge,
and which passes the frozen anti-wash-trading policy. Challenge work may earn a score under a
separate allowance, but it never increases $O_e$.

\subsubsection{Control-Domain Aggregation and Concentration}

Addresses are aggregated before scoring. For control domain $a$, define raw epoch score

\[
  S_{a,e}=\sum_{j\in J_e:\,a(j)=a} v_j .
\]

The published methodology then applies a concave tail and an absolute score cap:

\[
 \widehat S_{a,e}=
 \min\!\left\{C_e,
 \begin{cases}
   S_{a,e}, & S_{a,e}\le H_e,\\[0.25em]
   H_e+\left\lfloor \alpha_e(S_{a,e}-H_e)\right\rfloor,
     & S_{a,e}>H_e,
 \end{cases}\right\}
 \qquad 0\le\alpha_e\le1.
\]

$H_e$ is the full-credit threshold, $\alpha_e$ the published marginal factor above it and
$C_e$ the maximum credited score for one control domain. This is a cap on credited score, not
a claim that one operator can never receive a large fraction of a thin epoch. Actual top-$k$
reward shares and the Herfindahl index remain mandatory disclosures. Splitting addresses does
not raise the limit when beneficial ownership, operator keys, payer relationships, hosting,
network identity or other published common-control signals link them. Because hidden common
control cannot be proven from chain data alone in every case, the domain graph and each
classification change are reproducible scorer inputs and are challengeable with the score
root.

Let $\widehat S_e=\sum_a\widehat S_{a,e}$. The epoch Merkle tree commits leaves of the form
$(a,\widehat S_{a,e})$, sorted by identity, together with the frozen methodology and rate-card
hashes. Duplicate identities are invalid rather than silently merged.

\subsubsection{Demand-Coupled Epoch Pool}

The native epoch pool uses exact integer arithmetic. Let $K_n/K_d$ be the configured emission
coefficient, $F_e$ the cold-start floor, $U_e$ the published per-epoch schedule ceiling,
$C_{\mathrm{AI}}$ the cumulative AIPoW cap and $M_{<e}$ cumulative prior AIPoW creation. The
candidate pool and final mint are

\[
  P_e=\min\!\left(
       U_e,
       \max\!\left(F_e,
       \left\lfloor O_e\frac{K_n}{K_d}\right\rfloor\right)
      \right),
  \qquad
  M_e=\min\!\left(P_e,\max(0,C_{\mathrm{AI}}-M_{<e})\right).
\]

This is the arithmetic implemented by the feature-gated native mint path. A positive cold-start
floor is an explicit, measurable acquisition subsidy, not organic demand, and must be reported
separately. Governance may set it to zero. A randomized challenge allowance in the scoring
methodology is bounded by

\[
  S^{\mathrm{challenge}}_e \le
  F^{\mathrm{challenge}}_e+
  \left\lfloor O_e\frac{\chi_n}{\chi_d}\right\rfloor,
\]

where the first term is a published cold-start allowance and $\chi_n/\chi_d$ is the configured
challenge multiplier. Related-payer and challenge settlements remain excluded from $O_e$, so
they cannot recursively expand the demand-coupled pool. If no valid final score commitment
exists after the grace period, the epoch is skipped with zero creation. No shortfall is owed,
queued, extended or back-filled.

\subsubsection{Distribution and Reward Maturation}

For a valid epoch with $\widehat S_e>0$, domain $a$ is entitled to

\[
  R_{a,e}=\left\lfloor
    M_e\frac{\widehat S_{a,e}}{\widehat S_e}
  \right\rfloor,
  \qquad
  \sum_a R_{a,e}\le M_e.
\]

A claimant supplies a Merkle proof, but payment always goes to the committed earning identity,
never to the transaction sender. Integer remainder and unclaimed value stay unassigned in the
epoch distributor; no operator may redirect them. The distributor snapshots an immediate
fraction $b/10000$, a stream length $n$ and maturation-epoch duration $\tau$. If $t_0$ is the
claim time, the matured amount at time $t$ is

\[
  A_{a,e}(t)=I_{a,e}+
  \left\lfloor
   (R_{a,e}-I_{a,e})
   \frac{\min\!\left(n,
     \left\lfloor\frac{(\min(t,t_f)-t_0)_+}{\tau}\right\rfloor\right)}{n}
  \right\rfloor,
  \qquad
  I_{a,e}=\left\lfloor R_{a,e}\frac{b}{10^4}\right\rfloor .
\]

For an un-forfeited claim $t_f=+\infty$. On proven fraud, $t_f$ is the forfeiture time and
the unmatured remainder is permanently void; value already matured remains payable. Maturation
is delayed payment for completed work and a period in which fraud remedies remain meaningful;
it is not a return for passively holding or staking TOS.

\subsubsection{Commit, Challenge and Native Settlement}

The scoring and issuance pipeline separates expensive interpretation from deterministic
consensus:

\begin{enumerate}
  \item Public chain data yields settled-work rows. The interim feed records only settled Task
        Escrows and paid Service Actor requests; a production receipt adds capability class,
        units, explicit evidence level and price-cap inputs.
  \item Independently implemented scorers apply the same frozen methodology, rate card and
        common-control inputs and must reproduce $O_e$, every $\widehat S_{a,e}$,
        $\widehat S_e$ and the Merkle root byte for byte.
  \item A committer bonds the root, total score, organic settled value, methodology hash and
        rate-card hash. A symmetric bonded public challenge can present contrary evidence.
        An unreviewed challenge fails the root safe after its deadline; an absent or rejected
        root cannot halt later epochs.
  \item After the full challenge window, the masterchain verifies the audited contract code
        hash, canonical address, approved reviewer, final status, economic tuple and frozen
        configuration. Collators and validators call one shared pure integer derivation for
        $M_e$; disagreement makes the block invalid.
  \item The settlement actor deploys an audited per-epoch distributor. Claims are
        permissionless, Merkle-verified, pro rata, paid only to the committed identity and
        subject to the snapshotted maturation schedule.
\end{enumerate}

Consensus does not decide whether an answer is intelligent. It decides whether the approved,
challenge-complete score commitment and exact supply arithmetic authorize a bounded mint.
This is an optimistic-verification boundary: its security depends on data availability,
independent recomputation, a live challenger set and a reviewer that cannot be chosen by the
committer. Verification of untrusted ML is itself an active research area; replication,
proof-of-learning and cryptographic proof systems carry materially different cost and
soundness trade-offs.\footnote{For one primary-system discussion of these trade-offs, see
Gensyn, ``Verde: a verification system for machine learning over untrusted nodes,'' 2025:
\url{https://blog.gensyn.ai/verde-a-verification-system-for-machine-learning-over-untrusted-nodes/}.}

\subsubsection{Attack Economics and Launch Criterion}

For a fraud strategy with maximum gain $G$, probability $d$ of detection and successful
challenge, slash or forfeiture loss $L$, and direct attack cost $C$, a conservative bound is

\[
  \mathbb{E}[\Pi_{\mathrm{fraud}}]\le(1-d)G-dL-C.
\]

The economic launch criterion is not the existence of a bond; it is evidence from adversarial
tests that the right-hand side is negative for the largest plausible single-epoch strategies,
including wash settlement, payer-earner collusion, forged evidence, scorer equivocation,
reviewer capture, address splitting, duplicate receipts and withholding evidence until after
the challenge window. The project must publish assumed $G,d,L,C$, sensitivity ranges and the
largest profitable counterexample found. Where reliable bounds cannot be established, the
affected evidence multiplier, cold-start floor or emission coefficient must be reduced,
possibly to zero.

The implemented repository contains the interim settled-work data plane; score-commitment,
challenge, settlement and distributor actors; an independent Merkle implementation and claim
path; governance-gated ConfigParams 90--93; a cumulative native-supply cap; and a shared
collator/validator mint derivation. The \texttt{capAipow} feature is off in the genesis
template. Production activation additionally requires the full settlement-receipt schema, a
versioned scoring methodology and control-domain policy, two independently authored scorers
that reproduce long-running shadow epochs, data-availability and challenge tooling, review of
the reviewer/governance trust model, external contract and consensus audits, adversarial
economic testing, and counsel review of then-current applicable law. None of the community
allocation is created before those gates pass, and this section promises no completion date,
reward rate, adoption level or token value.

\subsection{Broad Distribution and Concentration Controls}

The protocol avoids allocating the initial supply to founders, investors or a central
treasury. Its validator distribution is intended to broaden through recurring elections and
open network participation. Initial controls include:

\begin{itemize}
  \item four equal-weight original validators rather than one bootstrap validator;
  \item equal and capped bootstrap funding to separately disclosed controlling wallets;
  \item recurring protocol elections that replace the zerostate validator set;
  \item a 10,000-TOS minimum and 10,000,000-TOS maximum submitted stake;
  \item a four-validator, 40,000-TOS aggregate protocol minimum;
  \item an initial maximum effective-stake factor of one, which keeps the minimum
        four-validator elected set equal in effective weight;
  \item a maximum of 400 validators and 100 masterchain validators;
  \item public disclosure of beneficial control, key custody, funding, hosting, network,
        geography and related-party relationships; and
  \item public concentration metrics for the largest one, four and ten validators and for
        known related parties.
\end{itemize}

Four validators are only the bootstrap minimum, not a decentralization claim. The
recommended public-launch participation target is 64 independent operators, with a
long-term target of at least 75 eligible validators.

A later increase in the effective-stake factor requires sustained independent
participation, concentration analysis, election simulation and public governance review. It
also has an arithmetic precondition that is not a matter of review. The Elector pays each
elected validator on $\min(s_i, f \cdot s_{\min})$, where $f$ is the effective-stake factor
and $s_{\min}$ is the smallest stake in the elected set, and it refunds the remainder. The
largest share of effective weight one entry can therefore hold in a set of $n$ validators is

\[
  \frac{f}{f + n - 1},
\]

and requiring that this stay below one third gives

\[
  f < \frac{n - 1}{2}.
\]

The bound applies to the smallest set the configuration permits rather than to the set that
happens to be running, because the Elector may install exactly the minimum in any round.
Raising the factor to three therefore requires a minimum validator set of at least eight,
and the minimum must be raised first: increasing the factor ahead of it would leave a window
in which the configuration admits a set where one entry holds half the weight. Repository
tooling enforces both the bound and the ordering, and refuses to emit a signable proposal
for a combination that violates them.

These controls reduce initial administrative allocation and single-entry dominance, but no
protocol can guarantee that market ownership or operational control will never concentrate.
Stake-proportional rewards may compound existing holdings, one controller may attempt to
operate several identities, and early rewards necessarily go to the validators then elected.
Broad ownership therefore also depends on new validator entry, service payments, voluntary
transfers and transparent monitoring. TOS communications must describe these residual risks
and must not claim guaranteed decentralization.

\subsection{Supply Accounting, Burns and Transparency}

Supply reporting distinguishes:

\begin{itemize}
  \item \textbf{gross native creation}: every TOS created in zerostate or by finalized block
        creation, including amounts later burned;
  \item \textbf{outstanding supply}: gross native creation minus provably burned TOS;
  \item \textbf{circulating supply}: outstanding TOS excluding published system reserves,
        locked stake and other consistently defined non-circulating balances; and
  \item \textbf{validator block creation}: cumulative ConfigParam 14 creation after genesis.
\end{itemize}

Burning unused bootstrap funds or penalties reduces outstanding supply but does not erase
historical gross creation. Public reporting must derive creation and burn totals from
finalized chain data, not solely from wall-clock estimates. Before launch the project must
publish every genesis balance, validator and controlling-wallet identity, relevant
configuration value, zerostate hash and bootstrap transaction cap. During operation it must
publish current reward parameters, cumulative creation, burns, circulating-supply
methodology, trailing block rates, target-date projections, validator concentration and
every monetary configuration proposal and vote.

\subsection{Metering and Charging Are Distinct}

An operation may be metered even when the caller pays zero at the moment of execution.
Settlement policies may include:

\begin{itemize}
  \item immediate native-token payment;
  \item escrowed maximum charge with refund;
  \item prepaid balance or subscription allowance;
  \item recipient or organization sponsorship;
  \item a refundable anti-spam or performance bond;
  \item off-chain batching or a payment channel; and
  \item deferred aggregate settlement.
\end{itemize}

Metering provides resource attribution, quota enforcement and auditable usage. Charging
allocates economic cost under a policy. Profiles define both, but they are not the same.
Not every operation requires one on-chain transaction.

\subsection{Service Economic Flows}

The architecture can use native TOS for:

\begin{itemize}
  \item escrowed task budgets and agent payouts;
  \item refunds after rejection, cancellation or timeout;
  \item Agentic Service Operation settlement under explicit quotes and maximum charges;
  \item prepaid communication, storage, compute, event and tool allowances;
  \item refundable introduction, reservation or performance bonds;
  \item registry bonds and future verifier staking;
  \item defined discovery, relay, verification or storage services; and
  \item gas for messages and persistent actor state.
\end{itemize}

Economic activity is designed to arise from completed operations and explicitly defined
infrastructure services, not from merely remaining online, sending self-generated traffic or
declaring hardware capacity. Token ownership does not grant an agent unlimited automation authority, prove
service quality or override terminal policy.

Service transaction volume, validator fees, service-provider revenue and token market value
are distinct concepts. This paper makes no forecast or promise regarding any of them.

\subsection{Development Funding and Progressive Decentralization}

A zero team, investor, foundation and treasury allocation reduces genesis insider overhang,
but it does not answer how core development, audits, SDK maintenance, ecosystem support and
go-to-market work are funded. Before production launch, the project must publish the legal and
operating entities involved, material sources and uses of funds, conflicts of interest, at
least a multi-year maintenance plan and the boundary between commercial products and protocol
governance. Any future grant, fee, equity, service-revenue or token-funded program must be
disclosed separately from protocol issuance.

Early product work requires a coherent team that can iterate quickly. Credible neutrality
requires that this dependency decrease over time. The decentralization plan should therefore
track validator diversity, independent client and SDK contributors, third-party operation
publishers, governance participation, documentation and release-key control. The objective is
not premature decentralization; it is a staged path by which the founding team becomes one
important contributor among several after product-market fit and operational safety exist.

\subsection{Governance, Communications and Participant Responsibilities}

Ordinary authenticated configuration governance retains the ability to change block rewards,
validator counts, stake limits, election timing and other configurable protocol parameters.
Every monetary proposal should disclose its proposer, old and proposed values, activation
conditions, expected supply impact and recorded votes. The ability to change configuration
is a material governance risk and means neither the 500-million validator-creation target
nor the five-billion total-supply target is a hard technical cap.

Public descriptions of TOS must be accurate, balanced and consistent with released code and
accepted network configuration. They must separate implemented functions from plans, avoid
predictions of token value or return, disclose material technical and economic risks, and
avoid presenting a policy target as a guarantee. No operator may describe validator
compensation as risk-free interest or a reward available merely for purchasing or holding
TOS.

The base protocol does not replace obligations that may apply to exchanges, custodians,
wallet operators, payment intermediaries, token distributors or commercial service
providers. Each participant is responsible for determining whether its activities require
authorization, customer and counterparty checks, transfer records, transaction monitoring,
tax reporting, sanctions controls, consumer disclosures, data protection or other
jurisdiction-specific measures. Service-layer implementations should expose the records and
policy boundaries needed for their operators to implement those controls without placing
private customer data in public consensus.

\subsection{Risk and Sustainability Disclosures}

TOS may lose some or all market value. Transferability, exchange availability and liquidity
are not guaranteed. Users face key-loss, custody, smart-contract, consensus, validator,
governance, software, cybersecurity, service-provider, market and changing-regulatory risks.
Network faults, concentrated stake, configuration changes or insufficient fee demand may
impair operation or validator incentives. TOS is not redeemable for official currency by the
protocol and is not presented as a deposit, insured balance or protected investment product.

Consensus uses stake-bonded validators rather than proof-of-work mining or a proof-of-work
token giver. Validator nodes, networking and supporting infrastructure still consume energy
and hardware resources. AI terminals and service workloads have separate resource impacts
that must not be attributed to consensus alone. Material energy and environmental claims
should be based on a published methodology and measured deployment data; no unverified claim
of zero impact is made.

This section describes protocol design and policy rather than providing legal, tax or
investment advice. Classification and participant obligations depend on the facts of an
activity and may change as the network, services and applicable requirements evolve.


\section{Implementation Status}

This whitepaper separates repository implementation from architectural direction.
``Implemented'' means a repository contains the relevant node, contract, tooling or test
foundation. It does not imply public commercial deployment, external audit, production
service availability or adoption. ``Partial'' means useful primitives or a profile-specific
implementation exist but the generic interoperable product surface does not. ``Planned''
identifies specification or product work that remains.

\begin{longtable}{L{0.34\textwidth}L{0.17\textwidth}L{0.38\textwidth}}
\toprule
Capability & Status & Scope \\
\midrule
\endhead
Native TVM execution and actor messages &
\implemented &
Node, contracts, cells, gas and value-bearing messages \\
Masterchain and shardchain protocol &
\implemented &
Validator and full-node protocol stack \\
ADNL, DHT, RLDP, QUIC and TOS Sites &
\implemented &
Core transport and networking primitives \\
TOS DNS and resolver chaining &
\implemented &
Naming primitives; not a complete public registrar product \\
Agent Wallet profiles &
\implemented &
Keys, policy, funding, activation, runtime binding and inspection \\
Agent Account &
\implemented &
Deterministic deployment, controller actions, limits and rotation \\
Contract message opcode pattern &
\implemented &
Typed contract ABI operations with opcode and query/correlation fields \\
Task Escrow &
\implemented &
Claim, accept, reject, result, review, dispute, resolve and settlement \\
Agent and task read APIs &
\implemented &
Authenticated reads, filtering, timeout and error taxonomy \\
Capability Registry &
\implemented &
Per-actor entry, metadata hashes, bond, activity and reputation \\
Service Actor contract &
\implemented &
Access policy, pricing, rate limits, calls, responses and revenue \\
Dispute and Proof Attestation foundations &
\implemented &
Compact resolution and evidence commitments \\
Open Intent and Agreement model &
\planned &
Canonical Intent, proposal, acceptance, amendment, cancellation, Agreement and recourse
schemas plus signature domains and conformance vectors \\
Agentic Service Operation model &
\partialstatus &
Profile-specific operations and common concepts exist; the generic normative envelope,
namespace and lifecycle are not frozen \\
Messenger operation profile &
\partialstatus &
Off-chain Messaging Plane foundations; production interoperability, economics and release
gates remain incomplete \\
Mailbox operation profile &
\partialstatus &
Capability-bound deposit/read/delete foundations and stored acknowledgements; generic
pricing, relay lease, bond and conformance remain incomplete \\
Open operation namespace and registry &
\planned &
Versioning, domain namespaces, immutable definitions, mirrored indexes and compatibility \\
Generic request, quote and receipt envelope &
\planned &
Canonical schemas, resource vectors, errors, SDKs and independent conformance \\
High-frequency service settlement &
\planned &
Prepaid balances, channels, receipt roots, batching, closure and dispute \\
Public \texttt{.tos} registrar product &
\planned &
Ownership lifecycle, SDK, deployment and operator tooling \\
ARD compatibility profile &
\planned &
Pinned external version, catalog mapping, publisher binding, media types and conformance \\
TOS ARD Registry &
\planned &
Independent REST search, ingestion, federation, provenance, visibility and bounded crawling \\
Provider and Agent SDKs &
\planned &
Discovery, quote, capability, invocation, metering, receipt verification and settlement \\
Managed model and tool provider profile &
\planned &
Approved services, bounded scheduling, metering, evidence and receipts \\
Storage and event profiles &
\planned &
Lease, retention, fan-out, replay and settlement semantics \\
Commerce, booking and human-service profiles &
\planned &
Domain state machines, deposits, cancellation, fulfillment and dispute \\
Physical AI terminal &
\planned &
Offline operation, safe updates, real-time priority, isolation and fleet management \\
AIPoW settled-work data plane &
\implemented &
Interim indexer feed for settled Task Escrows and paid Service Actor requests; full operation
and receipt fields remain planned \\
AIPoW commitment and distribution actors &
\implemented &
Bonded root challenge, epoch settlement, Merkle claims, maturation and forfeiture; not
production-activated \\
AIPoW native issuance path &
\implemented &
ConfigParams 90--93, cumulative cap and shared collator/validator derivation;
\texttt{capAipow} is off by default \\
Pooled-stake contract and tooling &
\partialstatus &
Application-layer arrangement, not a protocol feature; no general user entry point and no
economic effect at effective-stake factor one \\
AIPoW production scoring and activation &
\planned &
Frozen methodology, operation eligibility, control-domain inputs, independent scorer
reproduction, audits, shadow epochs and governance gate \\
Public-testnet hardening and external audit &
\planned &
Persistent deployment, operation-level load testing and independent review \\
\bottomrule
\end{longtable}

\section{Verification and Release Gates}

No Agentic Operation profile is production-ready merely because its happy path works. Release
suites must cover:

\begin{itemize}
  \item canonical operation identifiers, revisions, schema and metering hashes;
  \item signature domains, replay, expiry, idempotency and unknown critical extensions;
  \item capability scope, delegation depth, revocation, quota and cumulative budget;
  \item quote expiry, price-vector arithmetic, maximum-charge enforcement and refunds;
  \item operation cancellation, partial completion, restart and duplicate suppression;
  \item receipt signatures, usage conservation, evidence levels and settlement references;
  \item cross-provider conformance for success and normalized failure cases;
  \item contract authorization, escrow conservation, cancellation, refunds and disputes;
  \item Messenger encryption and session safety, device binding and recipient policy;
  \item Mailbox byte, retention, replica, read, delete and delivery-attempt limits;
  \item spam campaigns, introduction-bond abuse, reputation gaming and recipient false
        positives;
  \item event, subscription, retry, referral and broadcast amplification;
  \item malformed and oversized payloads, attachment expansion and decompression;
  \item resolver, descriptor, quote, payment, invocation, receipt and restart flows on a
        multi-node TOS network;
  \item pinned-version ARD publication, \texttt{POST /search}, publisher verification,
        provenance, withdrawal, nested catalog and bounded federation behavior;
  \item SSRF, DNS rebinding, redirect and federation loops, prompt injection, stale rollback,
        equivocation, poisoned ranking and private-catalog leakage;
  \item adapter crash, timeout, disconnect, malformed input, out-of-memory and failed-payment
        cleanup;
  \item bounded RAM, accelerator memory, disk, cache, watcher, connection, retry and journal
        behavior under sustained load;
  \item channel and batch sequence, closure, timeout, dispute and receipt-root verification;
  \item physical real-time priority, network partition, stale policy, reconnect and
        idempotent reconciliation;
  \item signed update verification, compatibility rejection, power-loss activation,
        rollback and staged fleet abort;
  \item actuator isolation, local safety override, fleet revocation and bounded command
        fan-out;
  \item AIPoW score arithmetic, operation eligibility, rate-card caps, duplicate rejection,
        common-control aggregation, deterministic Merkle roots and independent scorer equality;
  \item AIPoW commitment and challenge deadlines, reviewer failure, malformed roots,
        canonical-address and code-hash binding, skipped epochs and data availability; and
  \item adversarial AIPoW wash settlement, payer-provider collusion, Sybil splitting,
        evidence forgery, scorer equivocation, challenger inactivity and reviewer capture.
\end{itemize}

Release gates include a public threat model, protocol conformance vectors, reproducible
builds, system-service or container hardening, migration and rollback procedures, at least two
independent interoperable implementations for a profile, independent review of externally
exposed and settlement-critical components, and observed testnet operation.

\subsection{Compatibility Gate}

An operation revision is frozen only after:

\begin{enumerate}
  \item canonical schemas and hashes are published;
  \item success, failure, retry, cancellation and receipt vectors are stable;
  \item two implementations produce and verify equivalent envelopes;
  \item downgrade and unknown-extension behavior is tested; and
  \item the old revision has a declared deprecation and coexistence policy.
\end{enumerate}

A new implementation may add a provider, transport or optimization without changing the
operation revision. A semantic or economic incompatibility requires a new revision.

\subsection{AIPoW Shadow Gate}

Keep \texttt{capAipow} inactive while completing the full settlement-receipt schema, operation
eligibility table, frozen scoring methodology, rate card, common-control inputs and public
epoch dataset. Run at least two independently authored scorers over sustained testnet shadow
epochs; publish every root, disagreement, challenge, concentration metric, attack simulation
and parameter sensitivity. A failed or delayed gate creates no emission debt.

\section{Roadmap}

\subsection{Foundation: TOS Core}

Maintain consensus, native TVM execution, networking, wallet and reusable actor-contract
foundations as a stable dependency. Application workloads and operation-profile release
cycles remain outside the validator.

\subsection{Phase 0: Agentic Operation Specification}

Freeze the common Intent, Agreement and Operation architecture with requirements from a first-party reference application,
at least one external agent builder and multiple unrelated service providers:

\begin{itemize}
  \item canonical Intent, proposal, negotiation, Agreement, amendment, cancellation and
        expiry semantics;
  \item principal, role, authority, acceptance, evidence-policy and recourse bindings;
  \item operation identity, namespace and immutable revisioning;
  \item canonical request, quote, capability, usage, receipt and error envelopes;
  \item replay, expiry, idempotency, cancellation and restart rules;
  \item resource-vector and maximum-charge semantics;
  \item sponsorship, allowance, bond, refund and settlement models;
  \item receipt and evidence levels; and
  \item conformance vectors and reference verifier.
\end{itemize}

\subsection{Phase 1: Software and Compute Beachhead, Communication Validation}

Deliver a machine-checkable software-work profile and the paid model and tool profile as the
initial external-acceptance wedges:

\begin{itemize}
  \item software tasks whose agreed repository state, tests, artifacts and review evidence
        can be checked by independent implementations;
  \item model and tool operations with live quotes, token/time/tool metering, bounded
        execution and signed usage receipts;
  \item a first-party agent application that uses only the public Agent SDK and operation
        envelopes;
  \item unrelated providers competing on the same frozen operation revision; and
  \item public reporting of quote conversion, repeat payers, completion, unit economics,
        settlement cost and receipt verification without AIPoW rewards.
\end{itemize}

Then use Messenger and Mailbox operations to validate E2EE payload handling,
capability-bound deposit/read/delete, receipt levels, quotas and introduction-bond
experiments. Each released profile requires at least two independently authored
implementations or provider adapters.

\subsection{Phase 2: Open Provider and Agent SDKs}

Deliver:

\begin{itemize}
  \item operation-definition tooling and signed provider descriptors;
  \item Agent SDK for discovery, quote, authorization, invocation and receipt verification;
  \item Provider SDK for admission, metering, receipt signing and restart recovery;
  \item independently deployable ARD Registry and operation-aware ranking;
  \item public \texttt{.tos} registrar and endpoint-binding workflow; and
  \item deterministic multi-node conformance harness.
\end{itemize}

\subsection{Phase 3: Storage, Events and Commerce}

Add storage lease, event publish/subscribe, booking, commerce, delivery and human-service
profiles. Keep their domain state machines separate while reusing common identity,
capability, budget, receipt and settlement semantics.

\subsection{Phase 4: Physical Edge Intelligence}

Deliver owner-operated reference terminals with disconnected local execution, bounded offline
authority, tamper-evident reconciliation, signed A/B updates, real-time priority, physical-I/O
isolation, independent safety interlocks and staged fleet management.

\subsection{Phase 5: High-Frequency Economic Infrastructure}

Add prepaid balances, payment channels, receipt-root commitments, netted batch settlement,
production relays, congestion pricing, replication, verifier markets and policy-aware
orchestration after operation and cleanup invariants are stable.

\subsection{Cross-Cutting AIPoW Gate}

Keep AIPoW inactive until operation eligibility, receipt schemas, independent scorer
reproduction, data availability, attack economics, audits and governance review satisfy
Section~\ref{sec:aipow}. Paid communication or self-generated operation volume does not
automatically qualify for issuance. Technical readiness is not sufficient: activation also
requires repeat organic payers, non-subsidized provider retention, published protocol and
service unit economics, and a sustainable development-funding plan.

\section{Non-Goals}

TOS does not aim to:

\begin{itemize}
  \item make every Agentic Service Opcode a native TVM instruction;
  \item require every valid operation to charge a non-zero fee;
  \item put private communication payloads, model prompts or application execution into
        blockchain consensus;
  \item replace HTTP, QUIC, WebSocket, MCP, A2A, OpenAPI or ARD with a TOS-only stack;
  \item treat discovery as authorization, live admission, payment or physical-control
        authority;
  \item run private model inference directly in consensus;
  \item store credentials, large files, full model outputs or raw sensor streams on-chain;
  \item rent bare GPUs or expose raw accelerator devices;
  \item accept arbitrary consumer-supplied containers, programs or models in a managed
        terminal;
  \item expose a public shell, Docker socket or unrestricted host interface;
  \item put blockchain consensus in a hard real-time or safety-control loop;
  \item expose raw CAN, GPIO, serial, fieldbus or actuator access as a public operation;
  \item reward a provider merely for remaining online, sending self-generated traffic or
        declaring capacity;
  \item hard-code one model provider, relay, storage backend, oracle, verifier or proof
        system;
  \item make third-party indexers authoritative for funds, permissions or settlement;
  \item give an agent runtime unrestricted owner custody by default; or
  \item bundle unrelated execution engines into the production node.
\end{itemize}

\section{Conclusion}

TOS Network is being built as the \textbf{foundational infrastructure---The Open System---for
the Agentic Internet}. Its mission is \textbf{Connect Every AI Agent}.

The Internet connected computers through common protocols. The Agentic Internet requires an
additional public layer through which autonomous participants can identify one another,
discover capabilities, express goals, negotiate exact terms, act within delegated limits,
exchange evidence and settle across organizational and platform boundaries.

The primary TOS lifecycle is therefore:

\begin{center}
\texttt{Identity -> Discovery -> Intent -> Negotiation -> Agreement}\\
\texttt{-> Bounded Execution -> Evidence/Receipt -> Settlement}.
\end{center}

Intent is not obligation. Discovery is not authority. Model interpretation is not policy
approval. Agreement is the authenticated commitment that binds roles, scope, budget,
evidence and recourse. Agentic Service Operations are the typed execution primitives that
fulfill those commitments; they are an important part of TOS, not its final definition.

Private payloads and workloads remain where they belong: in encrypted sessions, provider
runtimes, local storage and owner-controlled terminals. Providers retain sovereignty over
admission, hardware, models, data, pricing and safety. TOS Core supplies shared actor state,
value, commitments, disputes and final settlement where independent parties need a
tamper-resistant source of truth.

The architecture separates three opcode layers that are often confused. TVM opcodes execute
inside the deterministic blockchain VM. Contract message opcodes identify actor-level state
transitions. Agentic Service Opcodes identify economically accountable external actions.
Their shared name reflects a portable instruction vocabulary, not a requirement that every
action run in one virtual machine or that the operation layer define the whole network.

Metering does not mean every trusted message or local action must incur a non-zero charge. It
means no caller can impose unlimited storage, compute, delivery, attention or retry costs
without identity, capability, quota, budget or bond. A Receipt proves only the claims defined
by its schema, signer and evidence policy. These boundaries enable open interoperability
without turning the Agentic Internet into an open spam or false-proof relay.

TOS Core provides the implemented blockchain and networking substrate. Agent and task actors,
capability and service foundations, profile-specific Messenger and Mailbox work, and the
feature-gated AIPoW path establish important pieces. The open Intent and Agreement schemas,
operation namespace, common envelopes, provider SDKs, cross-provider conformance,
high-frequency settlement and additional profiles remain work to specify, implement, test
and independently review.

Bittensor demonstrates a complementary path: open subnet markets can organize and reward the
production of machine intelligence and other digital commodities. TOS addresses the broader
cross-agent infrastructure that can connect such resources to agents, owners, applications
and institutions through explicit identity, authority, Agreement, evidence and settlement.

Architecture alone is not adoption. The proof must come from machine-checkable work, repeat
organic demand, third-party embedding, cross-provider completion and a measured advantage for
TOS coordination and settlement. Token issuance follows that proof; it does not create it.
AIPoW remains inactive until product, economic, security and decentralization gates are
independently visible.

\begin{center}
\textbf{The Internet connected computers. TOS connects autonomous agents.}
\end{center}
\section*{Copyright and License}

Copyright \textcopyright{} 2025--2026 TOS Network.

The source code and documentation in the TOS repository are distributed under the GNU General
Public License v3.0. Protocol behavior is defined by released code, schemas and accepted
network configuration; this whitepaper describes the architecture and direction and is not a
substitute for implementation-level review. It is not an offer of securities or a promise of
investment return, service adoption, protocol revenue or token appreciation.


\end{document}
